\section{Available measurements and how they can enter the frozen OGS model}
\label{sec:measurements}

This section is based on the measurement catalogue in
\path{../cda_knowledge_base/measurements}.  The purpose is not to change the
provided OGS model, but to define defensible observation and parameter mappings for
inversion.  In the current XML the unknowns solved by OGS are liquid pressure
\(\pp\), temperature \(\TT\), and displacement \(\uu\).  The secondary fields exported by
the TRM process include velocity, saturation, stress, strain, and swelling stress.
Measurements must therefore enter through one of three routes:

\begin{enumerate}[leftmargin=1.5em]
  \item \textbf{State observation:} compare a measurement-derived quantity with an OGS
  output field after an observation operator, e.g. \(S_\ell\mapsto\theta\mapsto\rho\)
  for ERT.
  \item \textbf{Boundary-condition reconstruction:} convert atmospheric or niche
  measurements, especially RH, into liquid-pressure boundary curves.
  \item \textbf{Parameter-field inference:} infer a heterogeneous and anisotropic
  intrinsic permeability tensor \(\kperm(\bm x)\), and possibly porosity or retention
  parameters, then regenerate the XML/mesh parameter input for each inversion sample.
\end{enumerate}

The strongest first target is the intrinsic permeability tensor.  The current model
uses the constant diagonal parameter
\[
\texttt{k\_i}=\begin{pmatrix}1.00\times10^{-19}&0\\0&0.40\times10^{-19}\end{pmatrix}
\,\mathrm{m^2}.
\]
For the inversion this should be interpreted as a reference tensor only.  A practical
field parameterisation is
\[
\kperm(\bm x)=
\mathbf R(\theta)^\top
\begin{pmatrix}
k_\parallel(\bm x)&0\\
0&k_\perp(\bm x)
\end{pmatrix}
\mathbf R(\theta),
\qquad
k_\parallel(\bm x)=r(\bm x)k_\perp(\bm x),
\]
with bedding-related orientation \(\theta\).  The measurement chapter below explains
which data constrain this tensor directly and which constrain it only indirectly
through saturation, pressure, or water-content dynamics.

\subsection{Measurement inventory by model relevance}

The collected measurement catalogue is outside this report folder at
\path{../cda_knowledge_base/measurements/}.  Table paths are relative to that
catalogue root.  A machine-checkable version of the same inventory is kept in
\path{inversion_workflow/observation_manifest.json}; its validator currently checks
28 local file/table/mesh conditions across the nine observation groups with zero
failures.  A second, model-facing layer is kept in
\path{inversion_workflow/processed_observations/}.  It contains normalized CSV tables
generated from the collected workbooks and ZIP archives by
\path{inversion_workflow/scripts/build_processed_observations.py}; every row keeps
source-file and sheet/member provenance so values can be traced back to the raw
emails or TeamBeam transfers.

\begin{longtable}{P{2.7cm} P{5.0cm} P{3.7cm} P{5.0cm}}
\caption{Measurement folders, physical quantity, and intended model-entry route.}
\label{tab:measurement-entry}\\
\toprule
\textbf{Measurement} & \textbf{Files copied in the catalogue} & \textbf{Physical quantity} & \textbf{Model-entry route}\\
\midrule
\endfirsthead
\caption[]{Measurement folders, physical quantity, and intended model-entry route (continued).}\\
\toprule
\textbf{Measurement} & \textbf{Files copied in the catalogue} & \textbf{Physical quantity} & \textbf{Model-entry route}\\
\midrule
\endhead
\bottomrule
\endlastfoot

ERT &
\path{ert/source_files}: open-niche VTK ZIP, raw \texttt{tx0/ohm},
electrode positions, resistivity--water-content sheet/PDF, ERT slides. &
Electrical resistivity or resistivity change; water content only after empirical
conversion. &
Observation operator from OGS saturation/water content to resistivity, preferably on
the ERT inversion mesh or a projected common mesh.\\

NMR &
\path{nmr/source_files}: weekly \texttt{.dat} files, seasonal NMR ZIP
and extracted files, figures, NMR 2025/2026 presentations, coordinates. &
Hydrogen/NMR-derived volumetric water content and relaxation-time information at
sparse positions. &
Point/line observation of water content after deciding whether the target is total
water, mobile water, or a corrected free-water content.\\

Permeability pulse tests &
\path{permeability_pulse_tests/source_files}: 2021/2024/2025
workbooks and characterization papers. &
Interval-scale hydraulic/gas-test permeability or transmissibility around borehole
sections. &
Direct constraint on effective directional components or averages of
\(\kperm(\bm x)\); not a direct cell-wise tensor value.\\

Taupe/TDR &
\path{taupe_tdr/source_files}: \texttt{Taupe\_WC.xlsx}, coordinates,
borehole figures, ISU/TD slides. &
Apparent relative dielectric permittivity / TDR water-content proxy in EDZ distance
bands. &
Baseline-normalized saturation/water-content trend diagnostic; absolute residuals
require confirmed Taupe unit/calibration.\\

Suction/RH &
\path{suction_relative_humidity/source_files}: \texttt{OT\_RH5--8.xlsx},
location figure, annual report, TD slides/minutes. &
Relative humidity and suction/water potential in boreholes and niches. &
Boundary-condition reconstruction through Kelvin-equation liquid pressure, and
validation of desaturation state through the retention curve.\\

Coord./layout &
\path{coordinates_geometry_layout/source_files}: measurement
coordinate workbooks, figures, Tecplot layout. &
Measurement locations, borehole labels, and model coordinates. &
Mandatory geometric mapping for all observation operators.\\

Projection inputs &
\path{model_projection_inputs/source_files}: projection script, projected VTU mesh,
project/XML variant, 2D model archives, timestep notes. &
Mesh and OGS-input workflow material, not a physical measurement. &
Concrete bridge for mesh-based parameter fields and for aligning measurement outputs
with OGS timesteps.\\

Bedding/geology &
\path{bedding_geology_structure/source_files}: bedding image, bedding
page, characterization papers. &
Bedding and fault/structure orientation, not a state variable. &
Orientation prior for anisotropic \(\kperm\); structural explanation for local ERT
and pulse-test anomalies.\\

Other HM monitoring &
\path{other_hm_monitoring/source_files}: TD minutes, modelling slides,
levelling slides, HERMES input note, Tecplot layout. &
Extensometer, mini-piezometer, crackmeter, laser scan, levelling, RH/T context. &
Structured layout, levelling summary, and qualitative pressure/deformation validation;
Geoscope and laser-scan raw exports still missing for hard residuals.\\
\end{longtable}

\subsection{Source-file content pass and traceability}

The catalogue has also been checked by a second pass through the copied files
themselves, not only through the email summaries.  The generated index
\path{../cda_knowledge_base/measurements/measurement_content_deep_dive.md} contains
101 mined fact rows across the nine measurement folders, and each measurement folder
has a local \path{DATA_CONTENT_SUMMARY.md} plus
\path{derived_files/content_deep_dive/content_summary.csv}.  The mirrored navigation
copy is under \path{../cda_knowledge_base/measurements_info}.  This pass mines the
workbook sheets, raw \texttt{.dat}/\texttt{.ohm}/\texttt{.tx0} files, ZIP member
catalogues, and generated PDF/PPTX text extracts.  Table~\ref{tab:source-content-pass}
records the main data facts that should be treated as the report-level source
checklist.

\begin{longtable}{P{2.7cm} P{6.1cm} P{6.6cm}}
\caption{Second-pass mined source-file facts used as a report traceability checklist.}
\label{tab:source-content-pass}\\
\toprule
\textbf{Folder} & \textbf{Concrete copied data content} & \textbf{Modelling implication}\\
\midrule
\endfirsthead
\caption[]{Second-pass mined source-file facts used as a report traceability checklist (continued).}\\
\toprule
\textbf{Folder} & \textbf{Concrete copied data content} & \textbf{Modelling implication}\\
\midrule
\endhead
\bottomrule
\endlastfoot

ERT &
15 copied source files; the TeamBeam ERT ZIP has 1,736 indexed members
(1,675 \texttt{.vtk} fields and 61 timestep/list files) in dated folders from
November 2019 to December 2024.  The Thunderbird raw example has a 6,128-line
\texttt{tx0} file, 2,656-line open \texttt{ohm} file, 2,524-line closed
\texttt{ohm} file, and 72-electrode open/closed text files. &
ERT can be audited down to raw inversion and electrode files, but it remains an
indirect log-resistivity state observation.  The ZIP is intentionally retained
compressed while every member is indexed.\\

NMR &
54 copied source/data files.  The weekly 4S table has 21 rows from
2021-10-12 to 2022-03-20 with water content 9.897--13.169 vol.\% and mean
12.004 vol.\%.  The weekly 4E table has 149 rows from 2021-10-06 to
2025-09-09 with water content 8.509--12.229 vol.\% and mean 10.592 vol.\%.
The seasonal archive contributes 22 campaign tables and 294 position rows across
Niches 3 and 4, with water content 2.766--22.810 vol.\%. &
NMR is the strongest local water-content evidence, but the source pass also makes
the bound/interlayer-water problem visible: these values are total NMR water-content
observations and cannot be equated mechanically to \(nS_\ell\).\\

Permeability pulse tests &
6 copied source files, including two permeability workbooks.  The workbook pass
indexes chart sheets, interpreted 2021/2024 sheets, and raw pressure-decay sheets
with thousands of non-empty cells; the processed observation layer extracts
204 interpreted interval rows and 6,687 pressure-decay rows. &
This is the direct parameter-observation stream for \(\kperm\), but the gas pulse
tests are scalar finite-interval observations with uncertainty, not cell-wise tensor
components and not liquid relative permeability.\\

Taupe/TDR &
10 copied source files.  \path{Taupe_WC.xlsx} has sheets A3, A4, A7, and A8,
each with 212 dated rows from 2019-12-01 to 2025-10-10 and EDZ-band columns.
The sheet values span approximately 9.87--11.00 in the copied workbook units over
the four sheets. &
Taupe/TDR has enough local table structure for trend diagnostics by sensor and EDZ
band, but the absolute unit/calibration gate remains open; the report should not
present these values as calibrated saturation.\\

Suction/RH &
10 copied source files.  The OT\_RH5--8 workbooks contain daily two-column RH
time series: RH5 has 1,042 data rows (94.081--95.993\%), RH6 has 1,064 rows
(93.323--94.988\%), RH7 has 1,071 rows (13.188--97.272\%), and RH8 has
1,070 rows (13.203--96.219\%). &
RH5/RH6 are the cleanest copied open-twin RH streams.  RH7/RH8 low values are
explicit source-level outliers or failure candidates, so RH-derived pressure
boundaries must include sensor screening and provenance.\\

Coord./layout &
8 copied source files.  The September coordinate workbook has 109 rows and 10
columns; the November coordinate workbook has 39 rows and 8 columns.  Both are
indexed with numeric coordinate ranges and linked to the copied figures and
\path{VisualisationCDA.dat}. &
All point, line, interval, and borehole observations must pass through this
geometry layer before they are compared with the current Niche-4 OGS mesh.\\

Bedding/geology &
4 copied source files, including the bedding image, bedding-page PDF, and local
characterisation-paper copies. &
These files provide the structural prior for permeability anisotropy and local
fault/EDZ interpretation, not a direct residual.\\

Other HM monitoring &
12 copied source files.  The content pass indexes five PDFs, three PPTX decks,
the Tecplot-style \path{VisualisationCDA.dat}, and response-note placeholders for
missing numeric exports and uncertainty information. &
The current report can use these files for qualitative HM context and levelling
summary evidence, but hard Geoscope, laser-scan, crackmeter, and related residuals
still require numeric exports and uncertainty definitions.\\

Projection inputs &
87 copied source files, including four ZIP archives, 57 XML files, 15 VTU files,
the projection script, project files, timestep slides, and the extracted report ZIP.
The source pass confirms the projection workflow references
\path{generate_projections.py}, \path{specifications.pptx}, and
\path{cd_a_open_niche_quad.prj}. &
This folder is the implementation bridge from measurements to OGS: it defines the
mesh fields, projection support, timestep convention, and exchanged model versions
without changing the frozen source model equations.\\
\end{longtable}

\subsection{ERT: resistivity, water content, and saturation}

The ERT archive contains a long open-niche time series: 1,736 files, including
1,675 VTK files and 61 monthly/timestep lists from November 2019 to December 2024.
The processed tables \path{ert_archive_inventory.csv} and \path{ert_timesteps.csv}
catalogue those ZIP members without unpacking the full archive.  They list 1,691
timestep-list entries and explicitly flag 16 entries without a matching VTK member,
so an observation operator should not assume a one-to-one date--VTK map without this
check.  The companion processed files \path{ert_nmr_resistivity_pairs.csv} and
\path{ert_kruschwitz2004_relation.csv} normalize the local water-content/resistivity
calibration workbook.
The raw ERT example recovered from Thunderbird contains one \texttt{tx0} file with
both niches and separate \texttt{ohm} inversion input files for the open and closed
niches.  Gesa's accompanying note states that the raw \texttt{tx0} file contains
open-niche measurements \#1--\#2592 and closed-niche measurements \#2593--\#5184.

ERT does not measure water content or saturation directly.  It measures an electrical
response that is inverted to apparent/effective resistivity.  Ziefle et al. use the
CD-A NMR/ERT crossplot to establish an empirical relation between electrical
resistivity and volumetric water content; the paper writes an Archie-type relation
\(\rho=c\rho_{\rm fluid}\theta^{-m}\) and then uses the fitted relationship for
comparison with simulated water content \cite[p.~4453, Eq.~4 and p.~4457,
Sect.~5.1]{WaterContentEDZ2024}.  This is consistent with Archie's original use of
formation factor and saturation-resistivity relations in clean porous media
\cite[pp.~54--58]{Archie1942}.  However, CD-A is claystone, not clean
sandstone: clay-bearing systems can carry current both through the pore fluid and
along mineral--water surfaces, with cation-exchange/surface-conductivity terms and
saturation transforms entering the conductivity model
\cite[pp.~23925--23927 and p.~23930, Eq.~17]{Revil1998ShalySands}.  The CD-A ERT
calibration must therefore remain local and empirical rather than a universal
saturation transform.

For modelling, the correct direction is:
\[
S_\ell^{\rm OGS}(\bm x,t)
\quad\rightarrow\quad
\theta_{\rm model}(\bm x,t)=n\,S_\ell(\bm x,t)
\quad\rightarrow\quad
\rho_{\rm pred}(\bm x,t)
\quad\rightarrow\quad
\mathcal O_{\rm ERT}[\rho_{\rm pred}],
\]
where \(\mathcal O_{\rm ERT}\) includes projection/interpolation to the ERT mesh and any
depth/ring averaging.  The inverse direction, converting ERT to ``measured
saturation'', is useful for plots but should not be treated as raw truth because
fractures, bedding-parallel conductive paths, electrode geometry, inversion smoothing,
and the different ERT/OGS meshes all enter the observation.  The September 2025
slides explicitly list these as open inverse-problem questions, especially filtering
unstable points and handling the mismatch between the ERT computational domain and
the OGS domain \cite[PDF pp.~3, 16--17]{CDAModellingSlides2025}.

The workflow now makes the calibration part of \(\mathcal O_{\rm ERT}\) explicit in
\path{inversion_workflow/scripts/build_ert_observation_operator.py}.  The script fits
the paired NMR/resistivity rows and the smooth workbook relation columns and writes
\path{ert_water_content_resistivity_operator.csv}.  The current default first-test
relation is the workbook \path{model_data2019} curve,
\[
\rho_{\Omega{\rm m}} = 1.108\,\theta^{-1.58},
\qquad \theta=nS_\ell ,
\]
over the calibrated water-content range \(0.01\le\theta\le0.18\).  The direct Niche-4
paired NMR/resistivity fit is retained as a diagnostic, but its scatter is larger
(0.202 log10 resistivity units), so it should not be mistaken for a completed ERT
forward model.

The spatial part of the ERT operator is now represented separately by
\path{build_ert_spatial_projection_lookup.py}.  It reads a reference ERT VTK from
\path{ERT_meas_Niche_open.zip}, records the raw ERT inversion mesh
(\(3{,}126\) points and \(5{,}966\) triangular cells), applies the explicit provisional
transform
\[
x_{\rm model}=x_{\rm ERT}, \qquad y_{\rm model}=z_{\rm ERT}+500\,{\rm m},
\]
and maps ERT cell centroids to the OGS bulk mesh.  The generated
\path{ert_spatial_projection_lookup.csv} has 4,676 ERT cells inside an OGS cell and
2,035 cells that are also inside the approximate 1.5 m central support mask noted in
the local slides \cite[PDF pp.~16--18]{CDAModellingSlides2025}.  This turns the ERT
mesh mismatch into a reproducible lookup instead of an undocumented manual step.
However, the transform and the exact ``35 cm in rock'' near-niche support mask still
need confirmation against the agreed ERT/FEM coordinate convention before ERT
log-resistivity residuals are activated.
The generated ERT semantics audit
\path{inversion_workflow/processed_observations/ert_measurement_semantics.md}
records the same gates in row-level form.  It keeps 1,691 ERT timestep rows, of
which 1,675 have matching VTK members and 16 are blocked by missing VTK files.  The
reference projection lookup has 5,966 ERT cells; 4,676 are inside an OGS cell,
2,194 are inside the approximate 1.5 m support screen, and 2,035 satisfy both
conditions.  These 2,035 rows are the current provisional ERT support under the
explicit transform/support assumptions.  The audit also records the recommended
residual as \(\log_{10}\rho_{\rm pred}-\log_{10}\rho_{\rm ERT}\) and explicitly
excludes use of ERT as measured saturation, measured water content, or measured
permeability.

The run-local diagnostic
\path{inversion_workflow/runs/direct_fit_observation_run/ert_resistivity_diagnostic.md}
now performs that comparison for the direct reference run.  It samples
\(\theta=nS_\ell\) from the OGS outputs on the 2,035 provisional ERT-support cells,
converts the sampled water content with the default Kruschwitz relation, and compares
the prediction with the nearest ERT inversion field within a 10-day tolerance.  The
diagnostic compares 162,800 cell-time rows across 80 OGS output times; 3 output times
are outside the time tolerance.  Its area-weighted residual
\(\log_{10}\rho_{\rm pred}-\log_{10}\rho_{\rm ERT}\) has MAE 0.254 and RMSE 0.300.
This closes the local OGS-output sampling step, but it is still a diagnostic rather
than a likelihood term because the coordinate transform, exact near-niche support
mask, and ERT inversion uncertainty/correlation model are not yet accepted.

The cross-run audit \path{build_ert_candidate_discrimination_audit.py} then applies
the same compact diagnostic to 66 executed OGS runs with output fields.  Each run
compares 162,800 ERT cell-time rows over 80 output times.  Across the current
candidate family the area-weighted ERT MAE spans 0.019635573360798686 log10 units.
The ERT-only best run is
\path{broad_continuous_001_001_length_0p023m_shift_1p004}, with MAE
0.25407209193077057; the best active-objective run is
\path{local_basis_sampler_002_basis_024_det_l_0p0075_s_1p000}, with ERT MAE
0.2541796182834413.  The combined-objective/ERT-MAE correlation is
0.8894403368395667, so candidates that fit the current active permeability+NMR
objective better generally do not improve the provisional ERT residual.  This is
candidate-screening evidence only; it does not remove the transform, support-mask, or
ERT uncertainty gates.

The support-mask part of that gate is now quantified locally by
\path{build_ert_support_sensitivity_audit.py}.  The audit reuses the row-level ERT
diagnostic on six representative executed runs and restricts the current provisional
1.5 m support to tighter radial caps and inner/outer annuli.  It evaluates nine
support variants.  The ERT-only run
\path{broad_continuous_001_001_length_0p023m_shift_1p004} remains the best mean
support-rank candidate, with mean rank 1.67 and worst rank 4.0, but the tightest or
outer support subsets can swap the winner to
\path{local_basis_sampler_003_basis_029_det_l_0p0100_s_1p000}.  This strengthens the
diagnostic record because support sensitivity is now explicit; it still does not
authorize ERT likelihood weights until the agreed tunnel-contour support and
uncertainty/correlation model are fixed.

ERT constrains permeability indirectly: permeability changes the simulated
desaturation/resaturation dynamics, which then affects predicted resistivity.  The
ERT data are therefore valuable for heterogeneous \(\kperm(\bm x)\), but only through a
forward observation model; they are not a direct permeability measurement.

\subsection{NMR: total water signal versus OGS liquid saturation}

The NMR folder contains two weekly monitoring tables and a seasonal campaign archive.
The weekly 4S table has 21 rows from October 2021 to March 2022; the weekly 4E table
has 149 rows from October 2021 to September 2025.  The values are reported as
volumetric water content with confidence information and relaxation-time columns.  The
seasonal ZIP contains 22 \texttt{.dat} data files and 22 figures, covering Niche 3 and
Niche 4 campaigns between 2019 and 2025.  The processed tables
\path{nmr_weekly.csv} and \path{nmr_seasonal_profiles.csv} combine these into 170
weekly rows and 294 seasonal campaign-position rows.  The weekly table flags the
February--April 2025 4E detuning caveat directly in the row-level \path{caveat}
column.

The CD-A paper treats NMR as the more direct water-content measurement used to anchor
the ERT interpretation; Fig.~6 shows the volumetric water content measured by
single-sided NMR at repeated surface positions, and Fig.~7 uses the NMR/ERT crossplot
for the empirical resistivity-water-content conversion
\cite[pp.~4451--4453]{WaterContentEDZ2024}.  That does not automatically mean that
NMR equals OGS saturation.  NMR observes hydrogen-bearing water with relaxation
behaviour controlled by pore geometry, surface relaxation, and clay effects.  The CD-A
NMR 2026 presentation in the local archive specifically flags interlayer/bound water
as a live interpretation issue; general NMR petrophysics literature also treats clay
effects and T2 distributions as a key limitation when connecting NMR to capillary
pressure and producible/mobile water
\cite[PDF pp.~23--24]{NMR2026Local}\cite[pp.~761--763]{Kleinberg1996NMR}.
An open-access dual-scale NMR modelling study makes the same point from the forward
modelling side: T2 distributions are used to infer pore size and permeability, but
their response is affected by internal magnetic-field gradients and unresolved clay
micro-porosity, so clay-rich regions require either explicit microstructure or
effective parameters rather than a direct saturation reading
\cite[pp.~453--455]{Cui2022DualScaleNMR}.

The defensible modelling map is therefore:
\[
\theta_{\rm NMR}^{\rm obs}
=\theta_{\rm mobile}+\theta_{\rm bound/interlayer}+\epsilon_{\rm NMR},
\qquad
\theta_{\rm mobile}^{\rm model}\approx nS_\ell .
\]
The generated bound-water sensitivity audit in
\path{inversion_workflow/processed_observations/nmr_bound_water_sensitivity.md}
quantifies how large this caveat is for the current model.  With fixed
\(n=0.105\), 315 of 464 NMR target rows exceed the maximum possible
model-side \(\theta=nS_\ell\) if interpreted as mobile water without correction.
Within the currently usable Niche-4 mapped subset, 162 of 287 rows exceed \(n\).
The required positive bound-water/free-water offset for usable rows has
p95 \(=0.0402\) and maximum \(=0.1231\); a tested uniform subtraction of 0.05
maximizes the simple physical-row count but still leaves seven non-physical usable
rows.  This means a single global subtraction is not a clean correction.

The cross-run NMR candidate bias audit
\path{build_nmr_candidate_bias_sensitivity_audit.py} tests the ranking consequence of
that caveat across 66 executed OGS runs with 192 active NMR rows each.  Under the
current raw absolute-\(\theta\) objective, the best run is
\path{local_basis_sampler_002_basis_024_det_l_0p0075_s_1p000}, with combined
objective 3156.353066948979.  If each NMR label is allowed a non-negative constant
bound/interlayer-water bias, or if only within-label anomalies are compared so that
constant offsets cancel, the best diagnostic run becomes
\path{broad_continuous_001_003_length_0p021m}, with diagnostic combined objective
505.614305828437.  The rank correlation between the current raw objective and the
per-label-bias diagnostic is 0.6351529067946979.  Therefore the present permeability
ranking is conditional on the unresolved NMR free-water interpretation, not a final
NMR-robust inversion result.

The provisional NMR objective decision in
\path{inversion_workflow/nmr_objective_decision.md} therefore recommends a
within-label trend/anomaly residual as the first final NMR likelihood candidate,
while the internal decision register records a current not-default policy for the
active report state.  Thus the raw-theta active objective is left unchanged, and the
trend/anomaly residual is used as explicit scenario/provisional-likelihood evidence
unless the modelling team reopens the default-objective decision.  The residual compares
\[
\left(\theta_{\rm model}-\overline{\theta_{\rm model}}_{\ell}\right)
-
\left(\theta_{\rm NMR}^{\rm obs}-\overline{\theta_{\rm NMR}^{\rm obs}}_{\ell}\right)
\]
within each observation-family/measurement-label group \(\ell\).  Constant
bound/interlayer-water and campaign offsets then cancel to first order, so the NMR
stream constrains water-content changes rather than absolute mobile water content.
This residual now has an executable assembler mode,
\path{--state-objective-mode} \path{nmr_within_label_trend_anomaly}, and the separate
\path{inversion_workflow/nmr_trend_anomaly_active_objective.md} ranking validates it
against the diagnostic audit to numerical roundoff.  In the 66 full active
direct-permeability plus NMR runs, it selects
\path{broad_continuous_001_003_length_0p021m} with objective 505.614305828437; the
raw-objective incumbent is rank 14 under this residual.
The final NMR residual-policy gate in
\path{inversion_workflow/nmr_final_residual_policy_gate.md} records the resulting
promotion boundary: no final NMR policy is selected yet, raw absolute theta remains
the current-report default with caveats, within-label trend/anomaly is the preferred
provisional final-policy candidate, and the follow-up planner recommends pausing new
OGS batches for that promoted-NMR mode because the best unevaluated direct advantage
is below the materiality threshold.
The matching acceptance-record template
\path{inversion_workflow/nmr_final_residual_policy_acceptance_record_template.md}
is the signoff guardrail for that choice.  It contains the four NMR residual
options, currently records zero of one required primary approvals, is not ready to
apply, records no actual decision, changes no active objective, promotes no field,
and recommends no new OGS batch.
The package explicitly rejects a single global NMR offset as a final correction
unless collaborators provide a physical justification, and keeps porosity fixed
because NMR water content alone cannot separate porosity, saturation, and bound
water.

For inversion, NMR should initially be used as a calibration/validation anchor for
near-surface water-content trends and for checking the ERT conversion.  If a free-water
correction becomes available from Stephan's lab interpretation, the NMR likelihood can
be tightened.  Without that correction, use a bias/error term or compare changes
\(\Delta\theta\) rather than absolute saturation.

Important local caveats from the email archive:

\begin{itemize}[leftmargin=1.5em]
  \item Columns 4--5 of the weekly files can be ignored for modelling.
  \item Technical failures cause gaps.
  \item February--April 2025 at 4E likely overestimates water content by about
  \(1\) vol.\% because the device was detuned to the wrong frequency.
  \item Seasonal data before autumn 2021 are floor/ceiling only; September 2020 is
  floor only; winter 2024 is missing.
\end{itemize}

\subsection{Permeability pulse tests: scalar interval data versus tensor field}

The permeability folder contains two main workbooks:
\path{2025-09-05_CD-A_Permeability.xlsx} and
\path{Permeability_CDA_all_2025.xlsx}.  The latter includes 2021 data from
BCDA\_24/BCDA\_25 and 2024 data from BCDA\_32/BCDA\_33.  The September slides describe
modified COMDRILL double-piston packer pulse tests with a 10 cm test interval,
nitrogen injection up to 1 bar, and pressure-decay monitoring
\cite[PDF p.~5]{CDAModellingSlides2025}.  The reported value is
an interval-scale intrinsic-permeability estimate for a 3D volume around the tested
borehole segment, with experimental/evaluation uncertainty of about half an order of
magnitude.  The processed files \path{permeability_interpreted_values.csv} and
\path{permeability_pressure_decay.csv} extract 204 interpreted interval rows and
6,687 raw pressure-decay rows while preserving source workbook, sheet, row, campaign
year, borehole block, depth, permeability, transmissibility, and pressure-time
provenance.

This is the most direct measurement class for \(\kperm\), but it still does not give a
cell-wise tensor.  It is an effective scalar or directional response of a finite test
volume and borehole orientation.  The CD-A characterization paper reports that pulse
tests indicate a stronger EDZ and permeability increase in the open twin than in the
closed twin, and later discusses interval permeability values around \(10^{-16}\) to
\(10^{-13}\,\mathrm{m^2}\) depending on niche, depth, and anomalies
\cite[PDF pp.~1, 11, Sect.~4.3 and Sect.~6]{Ziefle2024Characterization}.  The
workbooks and slides further distinguish 2021 (about 1.5 years after excavation) and
2024 (about 4.5 years after excavation).

The nitrogen-gas detail is crucial.  The OGS TRM Darcy law uses liquid intrinsic
permeability multiplied by liquid relative permeability:
\[
\vW=-\frac{\kperm k_{\rm rel}(S_\ell)}{\mu_w}\nabla p .
\]
Gas pulse tests can estimate the medium's intrinsic permeability only after
accounting for gas slippage and test conditions.  Klinkenberg showed that apparent gas
permeability depends on reciprocal mean pressure and that extrapolation to infinite
pressure gives the liquid-equivalent permeability constant
\cite[pp.~200, 205--207]{Klinkenberg1941}.  In CD-A, this means a nitrogen pulse-test
value should enter the inversion as a noisy effective observation of intrinsic
permeability, not as liquid hydraulic conductivity and not as \(k_{\rm rel}\kperm\).
This separation is also consistent with Opalinus Clay gas-transport literature:
gas flow in low-permeability claystone depends on intrinsic permeability, porosity,
gas pressure, water saturation, capillary threshold, and hydromechanical state; once
gas entry is exceeded, mobility is controlled by intrinsic permeability together with
relative-permeability and retention relations
\cite[pp.~124--125, Sects.~2.1--2.2]{Marschall2005}.  The pulse-test interpretation
therefore informs \(\kperm\), while liquid relative permeability and retention remain
separate OGS constitutive choices.

Recommended observation operator:
\[
d_i^{\rm pulse}
\approx
\frac{1}{|\Omega_i|}
\int_{\Omega_i}
\bm e_i^\top \kperm(\bm x)\bm e_i \, w_i(\bm x)\,d\bm x
\quad \text{with lognormal error,}
\]
where \(\Omega_i\) is the interval influence volume, \(\bm e_i\) is the dominant
borehole/test direction or sensitivity direction, and \(w_i\) is an approximate
sensitivity weight.  If the direction/sensitivity is not reconstructed, use the pulse
data as order-of-magnitude constraints on EDZ zones rather than point tensor entries.

The working folder now implements the first direct-target version of this idea in
\path{inversion_workflow/scripts/build_permeability_observation_targets.py}.  The
script reads the interpreted permeability rows, the borehole line samples, and the
mesh lookup table, then writes \path{permeability_observation_targets.csv} and
\path{permeability_observation_cells.csv}, plus
\path{permeability_missing_geometry_audit.md} for interpreted rows that have source
orientation evidence but no labelled endpoints.  It does not treat the scalar gas
pulse-test result as a tensor component.  Instead, it records a noisy interval-scale
log-permeability target with cell weights over a nominal 10 cm along-borehole
interval.  The current status is: 75 rows are usable as initial direct constraints in
the local OGS mesh, 27 rows map outside the current mesh and are retained but excluded
from the first fit, 98 older rows have interpreted values but lack matching borehole
endpoint geometry, and 4 rows have no positive interpreted permeability.  The 98
unmapped rows are grouped into five audit entries: BCD-A24 and BCD-A26 are retained
as vertical tests, BCD-A25 and BCD-A27 as horizontal tests, and BFM-D19 as nearly
horizontal evapometer data, but none is activated until labelled endpoint geometry or
an accepted digitized trace can project the interval into OGS cells.  The companion
request package
\path{inversion_workflow/processed_observations/permeability_endpoint_geometry_request.md}
turns this into a concrete follow-up item: it asks for labelled start/end
coordinates, coordinate frame, depth-zero reference, along-borehole direction, and
interval-position convention for BCD-A24, BCD-A25, BCD-A26, BCD-A27, and BFM-D19.
The same builder writes a 98-row blocked-observation table, so the inactive values
can be reactivated reproducibly if the endpoints or an approved digitized trace are
provided.
The generated semantics audit
\path{inversion_workflow/processed_observations/permeability_measurement_semantics.md}
keeps this decision machine-checkable: 204 interpreted rows are present, 200 have
positive permeability, 75 are active direct candidates, 27 are outside the current
mesh, and 98 are blocked only by missing endpoint geometry.  The active subset comes
from BCD-A32 and BCD-A33, with usable log10 permeability values spanning
\(-21.0\) to \(-12.10\)\,\(\mathrm{m^2}\).  The audit explicitly labels every row as
a nitrogen pulse-decay scalar interval observation of intrinsic permeability and
records that it is not a measurement of hydraulic conductivity, saturation, liquid
relative permeability, or a direct tensor component.

The companion evaluator
\path{inversion_workflow/scripts/evaluate_permeability_targets.py} is the first
machine-checkable likelihood component.  It reads a generated mesh field
\path{k_i_rd}, evaluates \( \bm e^T\kperm\bm e \) on the selected interval cells, and
compares predicted and observed permeability in log10 space.  For the current
\path{smoke_test} field, duplicate-aware weighting leaves an effective 52 direct
targets, with a weighted RMSE of 2.70 log10 units.  The best single global multiplier
is only \(10^{1.2285}\approx16.9\), reducing the RMSE to 2.40 log10 units.  This
baseline confirms that the current generated field is only a prior/proposal field;
the pulse-test data require heterogeneous local high-permeability intervals rather
than a uniform rescaling.

A first diagnostic direct fit is implemented in
\path{inversion_workflow/scripts/fit_permeability_field_from_targets.py}.  It scales
only the mesh cells selected by usable pulse-test intervals, preserves the existing
tensor orientation and anisotropy, and stores the applied log10 multiplier as mesh-cell
diagnostic fields.  In \path{runs/direct_permeability_fit/}, this reduces the
duplicate-aware direct-target RMSE from 2.70 to 1.61 log10 units by adjusting 30 cells,
with applied log10 multipliers from \(-1.48\) to \(+5.80\).  This is useful for
locating forced high-permeability intervals, but it is not yet a full OGS inversion
because it does not rerun the THM problem or include ERT, NMR, Taupe, or RH state
observations.

The prior/proposal space has also been tested explicitly with
\path{run_direct_permeability_prior_sweep.py}.  A 12-candidate sweep generated
heterogeneous anisotropic fields using the two usable pulse-test segment orientations
and the bedding-prior angle as tensor-angle candidates.  The best generated prior
field has RMSE 2.41 log10 units and objective 605.00, with an optimal uniform
multiplier of only 1.15; after that multiplier the RMSE remains 2.41.  This improves
on the original smoke-test prior but remains far worse than the direct cell-fit
diagnostic.  The conclusion is that the measured pulse-test intervals require local
permeability structure, not only a different global scale or a generic random
anisotropic prior.

A smoother interval-anchored field is implemented in
\path{fit_smooth_permeability_field_from_targets.py}.  This script takes the same
duplicate-aware target-derived log10 multipliers used by the direct cell diagnostic,
but spreads them from the 30 anchor cells through a finite-cutoff Gaussian kernel.  It
therefore keeps the tensor orientation and anisotropy while producing a neighbourhood
of affected cells instead of isolated point corrections.  The tested length scales
show the expected tradeoff: 0.05 m affects 471 cells and gives RMSE 1.67 log10 units
with objective 291.29, while 0.80 m affects 8,354 cells and relaxes to RMSE 2.13 with
objective 470.95.  An extended search now also varies shift damping.  Its best direct
candidate uses a 0.025 m smoothing length, affects 208 cells, and reaches RMSE 1.61
log10 units with objective 269.97.  This is still smoother than the 30-cell direct
diagnostic, but more localized than the earlier 0.05 m proposal.  The candidate has
been passed through the per-candidate driver as
\path{candidate_smooth_0p025m_search_driver}; it is currently the best direct-fit
field before state-observation comparison.

Because the direct objective alone favours the strongest local correction, the
workflow now also ranks the smooth candidates with
\path{rank_regularized_permeability_candidates.py}.  The ranking is not a calibrated
geological prior; it is a transparent selection aid that separates direct pulse-test
misfit from the sum of squared cell-wise log10 permeability updates.  The current
Pareto tradeoff set contains three 0.025 m support fields: the full-shift
\path{length_0p025m} field (objective 269.97, RMSE 1.61), the 0.75-shift field
(objective 301.16, RMSE 1.70), and the 0.5-shift field (objective 392.72, RMSE
1.94).  A weak field-update penalty still selects the full-shift field, a moderate
penalty selects the 0.75-shift field, and a strong penalty selects the 0.5-shift
field.  At least the direct winner and one damped regularized candidate should be
carried into the first OGS state-output comparison, so that the eventual inversion is
not judged only on direct permeability overfit.

That handoff is now reproducible with
\path{run_regularized_ogs_candidate_set.py}.  The script reads the Pareto/scenario
winners and prepares all three fields through the same candidate harness.  The
current executed set runs \path{length_0p025m},
\path{length_0p025m_shift_0p750}, and \path{length_0p025m_shift_0p500} through OGS
with the collected 6.5.4 SIF via Dockerized Apptainer.  Each run produced 83 output
VTUs, 37,184 sampled state rows, and 192 active sampled NMR objective rows.  The
combined objectives are 3166.61, 3198.49, and 3289.65, respectively.  The state
objective is similar for the three fields (about 2897), so this first combined
ranking is still driven mainly by the direct permeability term; it is an OGS-backed
candidate comparison, not yet a final inversion sampler.

The next-batch decision is now recorded by
\path{build_adaptive_combined_candidate_plan.py}.  The first adaptive batch has now
executed \path{length_0p050m}, \path{length_0p050m_shift_0p750}, and
\path{length_0p075m} through the same OGS harness.  Each run produced 83 output VTUs,
37,184 sampled state rows, and 192 active NMR objective rows.  Their combined
objectives are 3276.57, 3300.91, and 3353.88; none improves on the
\path{length_0p025m} objective of 3166.61.  The broader-support batch also increases
the sampled NMR state objective, so the executed state term is no longer flat over
that six-candidate set.  A subsequent local-refinement grid around the 0.025 m support
field executed \path{length_0p013m} and \path{length_0p019m}; both keep the same
direct RMSE of 1.61 log10 units and improve the combined objective to 3159.66 and
3164.96, respectively.  The first local-bracketing batch then executed
\path{length_0p013m_shift_0p875}, \path{length_0p013m_shift_1p125}, and
\path{length_0p031m}.  Their combined objectives are 3167.34, 3167.21, and 3219.35,
so neither shift perturbation nor the 0.031 m full-shift support improves on the
0.013 m local best.  A finite-candidate Bayesian proposal then executed
\path{length_0p019m_shift_1p125}, \path{length_0p019m_shift_0p875}, and
\path{length_0p025m_shift_1p125}; their combined objectives are 3172.69, 3172.37,
and 3173.39, again worse than the 0.013 m local best.  A first continuous-proposal
batch then executed \path{length_0p006m_shift_0p992}, \path{length_0p007m}, and
\path{length_0p007m_shift_0p972}; their combined objectives are 3156.78, 3157.74,
and 3158.16.  A focused lower-support continuous batch then executed
\path{length_0p004m_shift_0p995}, \path{length_0p003m_shift_0p986}, and
\path{length_0p004m_shift_0p994}; their combined objectives are 3156.38, 3156.47,
and 3156.38.  A repeatable loop driver,
\path{run_continuous_inversion_loop.py}, then refreshed the lower-support proposal,
executed \path{length_0p003m_shift_1p006},
\path{length_0p005m_shift_0p985}, and \path{length_0p004m_shift_1p014}, and merged
those rows into a cumulative lower-support evidence table.  Their combined objectives
are 3156.37, 3156.49, and 3156.44.  A second loop then executed
\path{length_0p004m_shift_0p992}, \path{length_0p004m_shift_1p020}, and
\path{length_0p006m_shift_0p975}; their combined objectives are 3156.40, 3156.53,
and 3156.70, so none improves the first-loop incumbent.  The current best executed
smooth field is therefore still \path{length_0p003m_shift_1p006}.  A broad
continuous evidence batch then executed \path{length_0p023m_shift_1p004},
\path{length_0p022m_shift_0p998}, and \path{length_0p021m}; their combined
objectives are 3159.06, 3169.30, and 3169.31, so the broad uncertainty-ranked batch
also did not improve the incumbent.  The generalized loop driver then ran a repeatable
broad iteration over \path{length_0p015m}, \path{length_0p016m_shift_0p968}, and
\path{length_0p010m}; their combined objectives are 3161.36, 3165.33, and 3157.70,
again leaving the incumbent unchanged.  The regenerated adaptive planner treats 32
executed smooth-family candidates as evidence and proposes the next
local batch led by \path{length_0p008m_shift_1p011},
\path{length_0p007m_shift_1p020}, and \path{length_0p007m_shift_1p022}.  This is
evidence that the lower-support continuous region has been useful, but the focused
local family is now near a plateau and still not a substitute for a production
optimizer or ensemble sampler.

The separate anisotropy sensitivity pass
\path{build_anisotropy_sensitivity_plan.py} checks a different question: whether the
direct pulse-test layer prefers a different global principal direction or anisotropy
ratio than the current bedding-prior tensor.  Starting from
\path{length_0p003m_shift_1p006}, it preserves each cell's geometric-mean
permeability and sweeps 35 angle/ratio combinations.  The best direct candidate is
the unmodified \(144^\circ\), 2.5:1 field itself, with direct objective 269.83549 and
zero improvement over the baseline; the next candidate, \(140^\circ\) with the same
ratio, is already worse by 0.05086 objective units.  This does not prove that
anisotropy is known perfectly, but it means a global angle/ratio-only OGS batch is
not a good use of computation.  The later local-basis, local-anisotropy,
cross-stream hybrid, structural/EDZ, and support lower-bound audits show that the
next permeability search should wait for a support/likelihood decision or for
currently gated streams such as ERT/Taupe/RH before adding broader tensor degrees of
freedom.

The first richer local-basis diagnostic is recorded by
\path{build_local_basis_sampler_plan.py}.  It starts from the same
\path{length_0p003m_shift_1p006} incumbent, keeps the governing OGS model frozen, and
builds 30 residual-derived anchor cells from the active permeability pulse-test
targets.  The sampler scores 131 direct pulse-test candidates in which local
coefficients can move independently around those anchors while only the run-local
\path{k_i_rd} tensor magnitude field changes.  The best direct row,
\path{basis_004_det_l_0p0015_s_1p000}, improves the direct pulse-test objective from
269.83549 to 269.81806, with weighted RMSE changing from 1.610768 to 1.610715 log10
units and applied log10 increments between \(-0.03472\) and \(0.00885\).  This is a
true but very small direct-layer improvement.  The selected OGS batch contains three
diverse local-basis candidates affecting 30, 37, and 47 cells.  The batch has now
been executed through the release-gated OGS harness.  All three runs returned code 0,
passed the release gate, sampled 192 active NMR state rows, and improved the combined
incumbent slightly: \path{basis_024_det_l_0p0075_s_1p000} has combined objective
3156.35307, compared with 3156.37301 for the previous
\path{length_0p003m_shift_1p006} incumbent.  The small improvement makes it the
current accepted best field, but it is still a local deterministic proposal result,
not a converged inversion.

The local tensor-anisotropy screen then checks whether the next field family should
release orientation or anisotropy ratio locally instead of continuing magnitude-only
updates.  The script
\path{build_local_anisotropy_sampler_plan.py} starts from the accepted local-basis
incumbent \path{basis_024_det_l_0p0075_s_1p000}, preserves each cell's
geometric-mean permeability, and changes only the local principal direction and
anisotropic split around the 30 mapped pulse-test anchor cells.  It scores 80 direct
pulse-test candidates across residual-sign, segment-parallel, segment-perpendicular,
and local-isotropization modes.  The best candidate,
\path{local_anis_isotropize_l0p015_s0p5_r1}, is worse than the incumbent: direct
objective 270.00798 versus 269.81806, a \(+0.18992\) objective increase.  Thus the
active direct permeability data do not justify a local tensor-orientation or
anisotropy-ratio OGS batch at this stage.

The cross-family production convergence record is now recorded by
\path{build_production_sampler_convergence_audit.py} and
\path{runs/production_sampler_convergence/PRODUCTION_SAMPLER_CONVERGENCE.md}.  This
audit standardizes 54 accepted smooth-family OGS rows and 11 accepted
local-basis rows, writes a stage-by-stage convergence history, and scores a
394-candidate smooth/local-basis pool with a random-forest surrogate for the sampled
NMR state objective plus the known direct permeability objective.  Five production
handoff rounds have executed thirty candidates; all passed release gates, sampled
192 active NMR rows, and were worse than the incumbent.  The best production row,
\path{basis_023_det_l_0p0075_s_0p750}, has combined objective 3156.35476; the fifth
round's best row, \path{length_0p004m_shift_1p031}, has combined objective
3156.78544, while the accepted local-basis incumbent remains
\path{basis_024_det_l_0p0075_s_1p000} at 3156.35307.  Of the 329 remaining
unexecuted diagnostic rows, 210 have materialized mesh files and are eligible for a
release-gated execution handoff; 119 random local-basis diagnostics are kept in the
score table but excluded from the batch because they do not yet have mesh files.  The
refreshed top executable handoff row is \path{length_0p007m_shift_1p020}, with
lower-confidence-bound objective 3156.94007 and a diagnostic probability of
improvement of 0.0166.  Because the local-basis improvement over the previous smooth
incumbent is only 0.01994 objective units and the cross-validated state-objective RMSE
of the surrogate is about 10.43, any next batch should be treated as convergence
evidence, not as proof that the inversion is complete.  The generated
\path{PRODUCTION_SAMPLER_DECISION.md} therefore recommends
\path{pause_active_production_sampling} for the current direct-permeability plus NMR
objective: all five production rounds failed to improve the incumbent, the next top
handoff probability of improvement is low, and all proposed lower-confidence bounds
are above the incumbent.  This pause is scoped only to the current active objective;
it does not close the ERT, Taupe/TDR, RH, or later-parameter gates.

The older smooth-family Bayesian proposal layer also writes
\path{runs/bayesian_candidate_proposal/OPTIMIZER_CANDIDATE_PROPOSAL.md}.  This is a
Bayesian surrogate over the already generated finite-grid and continuous smooth-field
pool: it records the 32 executed smooth-family candidates as evidence, keeps all 32 represented in
the generated-candidate score table, reconstructs fallback feature rows from executed
search results where proposal directories had been overwritten, adds the known direct
permeability objective for each unexecuted mesh, and ranks by a
lower-confidence-bound acquisition.  Its current top proposal is
\path{length_0p003m_shift_0p972}; the previous broad top proposals have now been
executed and proved worse than the incumbent.  The artifact is best read as a
transparent proposal diagnostic; a production inversion still needs an optimizer or
ensemble sampler loop over accepted fields.
The first continuous proposal layer is now recorded in
\path{runs/continuous_bayesian_candidate_plan/CONTINUOUS_CANDIDATE_PLAN.md}.  It
samples support length and shift scale outside the fixed finite grid, writes 136 new
run-local permeability meshes, evaluates the direct pulse-test objective for each,
and uses the current generated-table evidence intersection to predict the NMR
state-objective component.  Its
current proposed OGS batch, excluding already executed continuous rows, is
\path{length_0p009m_shift_0p979}, \path{length_0p008m_shift_1p040}, and
\path{length_0p010m_shift_1p050}.  Those rows are not yet accepted evidence; they must be
executed and sampled through the same release-gated OGS harness before comparing them
with the current \path{length_0p003m_shift_1p006} best field.  The focused
lower-support continuous plan separately proposes
\path{length_0p003m_shift_0p972}, \path{length_0p004m_shift_1p031}, and
\path{length_0p006m_shift_0p970}; after two executed lower-support loops its best
probability of improvement is effectively zero, so this local plan is now mainly a
plateau diagnostic.

\subsection{Taupe/TDR: dielectric proxy and EDZ-band trends}

The Taupe folder contains \texttt{Taupe\_WC.xlsx} with sheets A3, A4, A7, and A8.  Each
sheet has 212 rows from 2019-12-01 to 2025-10-10 and columns for EDZ distance bands
\texttt{0-50}, \texttt{0-10}, \texttt{10-20}, \texttt{20-30}, \texttt{30-40}, and
\texttt{40-50}.  The processed \path{taupe_tdr_bands.csv} table reshapes this workbook
to 5,088 sensor--date--EDZ-band rows and keeps the workbook column name and EDZ-band
limits in centimetres.  The 2026 ISU slides identify the method as differential TDR
using TAUPE sensor cables \cite[PDF p.~1]{TaupeISU2026Local}.  The CD-A
characterization paper lists Taupe as measuring
apparent relative dielectric permittivity and classifies it as saturation-related,
with limited information for EDZ/fault-zone/water-content interpretation
\cite[PDF pp.~4--5, Tables~1--2]{Ziefle2024Characterization}.

TDR methods infer water content because liquid water has a high dielectric
permittivity relative to mineral grains and air; the classic Topp et al. relationship
is a soil-calibration result, not an Opalinus-specific law
\cite[pp.~574--577]{Topp1980TDR}.  An accessible TDR review makes the same modelling
caution explicit from the measurement side: TDR estimates water content through bulk
permittivity, but porous media contain matrix, air, liquid water, and sometimes
bound water, and the measured response can also include electrical-conductivity and
dispersion effects \cite[pp.~444--446]{Robinson2003TDRReview}.  The local TD slides
use a simple volumetric mixing
expression
\[
\epsilon_r=\epsilon_{\rm rock}(1-\phi)+\epsilon_w\phi S_w,
\qquad \epsilon_w\approx 80,
\]
which is useful as a first-order link to saturation but should not be treated as a
fully calibrated Taupe forward model \cite[PDF p.~8]{CDAModellingSlides2025}.

The current workflow therefore writes an explicit Taupe/TDR trend-operator artifact
rather than an absolute saturation likelihood.  The script
\path{build_taupe_observation_operator.py} reads the 5,088 Taupe rows, the Taupe
line-sample mapping, and the OGS mesh porosity field \path{n_rd}.  It writes
\path{taupe_tdr_trend_operator.csv}, \path{taupe_tdr_series_summary.csv},
\path{taupe_tdr_observation_operator_summary.json}, and
\path{taupe_tdr_observation_operator.md}.  For each of the 24 sensor--EDZ-band
series, the first three finite workbook values define a baseline; later observations
are stored as raw differences, relative changes, and robust standardized anomalies.
Only A3 and A4 are mapped into the current local OGS mesh support, giving 2,544
trend-ready rows; A7 and A8 are retained but flagged as outside the present mesh
sampling support.  The intended model-side diagnostic is the corresponding
baseline-normalized trend in band-averaged
\(\theta_{\rm model}=\phi S_\ell\), not the raw Taupe number itself.

The follow-up semantics audit \path{build_taupe_semantics_audit.py} writes
\path{taupe_tdr_semantics_row_audit.csv},
\path{taupe_tdr_semantics_series_audit.csv},
\path{taupe_tdr_semantics_group_summary.csv},
\path{taupe_tdr_semantics_summary.json}, and \path{taupe_tdr_semantics.md}.  It
classifies the 2,544 A3/A4 rows as mapped trend-diagnostic candidates for runs with
sampled OGS outputs, while still blocking absolute calibration, and the 2,544 A7/A8
rows as outside the current mesh support.
The corresponding state-target table therefore contains 5,088 Taupe/TDR rows, but
none is active as an absolute state residual in the current candidate objective.

The run-local diagnostic \path{evaluate_taupe_tdr_trend_diagnostic.py} closes the
next comparison layer without changing the objective.  It matches the mapped A3/A4
Taupe rows to sampled OGS output times and forms the same-series standardized
residual
\[
  r_{\rm Taupe,trend} =
  \widehat{\Delta\theta_{\rm model}} -
  \widehat{\Delta T_{\rm Taupe}},
\]
where hats denote within-series robust standardization rather than a calibrated
absolute unit conversion.  In the direct reference run and the current local-basis
incumbent, the diagnostic compares 1,860 rows across 12 A3/A4 sensor--band series;
684 additional mapped A3/A4 rows are outside the available OGS output time horizon,
and all 2,544 A7/A8 rows remain outside the present mesh support.  The direct-run
standardized trend mean absolute residual is 1.8632.  This number is a screening
diagnostic for candidate behaviour, not a likelihood term, because Taupe units and a
sensor/band uncertainty model remain unconfirmed.

The cross-run audit \path{build_taupe_candidate_discrimination_audit.py} reuses the
same diagnostic across the executed run directories that already contain
\path{ogs_state_samples.csv}.  It currently evaluates 74 runs, including 66 runs with
the full active combined objective.  Across this executed candidate family the Taupe
standardized-trend MAE spans only 0.03687076560014302.  The best Taupe-only run is
\path{adaptive_combined_001_length_0p050m}, with MAE 1.829884354078035; the best
active-objective run is \path{local_basis_sampler_002_basis_024_det_l_0p0075_s_1p000},
with Taupe MAE 1.863211058631399.  This small spread means the present candidate
family is a weak discriminator for mapped Taupe trends, even before the
unit/calibration and weighting gates are solved.  The audit outputs are stored in
\path{taupe_candidate_discrimination_audit.csv},
\path{taupe_candidate_discrimination_series.csv},
\path{taupe_candidate_discrimination_summary.json}, and
\path{taupe_candidate_discrimination.md}, with catalogue copies under the Taupe/TDR
measurement folder.

The grouped-weight part of the Taupe gate is now tested by
\path{build_taupe_series_weight_sensitivity_audit.py}.  The audit uses the
series-level cross-run diagnostic for the 12 compared A3/A4 series over the 66 full
active-objective runs.  Because every compared series currently contributes 155 rows,
aggregate row weighting, equal-series weighting, equal-sensor weighting, and
equal-EDZ-band weighting all keep
\path{adaptive_combined_001_length_0p050m} as the best Taupe trend-screen run.  The
sensor-specific rankings are not identical: A3-only weighting prefers
\path{local_bracketed_003_length_0p031m}, A4-only weighting prefers
\path{adaptive_combined_002_length_0p050m_shift_0p750}, and the 12 individual
series have eight distinct best runs.  Thus the local evidence does not support a
single hard Taupe residual weight yet; sensor-specific uncertainty and grouped
weights remain part of the calibration gate.  A7/A8 also remain outside the current
local mesh support.

The same operator table records candidate absolute conversions for audit only.  If
the Taupe value is interpreted as volumetric water content in percent, 2,120 of 5,088
rows give mapped-band saturations in the physical interval \([0,1]\).  If the value is
interpreted as a Topp-style dielectric permittivity, none of the rows gives physical
mapped-band saturation for the current porosity field, which is a warning against
using that generic soil calibration as an Opalinus Clay law.  The local linear mixing
expression with \(\epsilon_{\rm rock}=6\) gives physical mapped-band saturations for
2,544 rows.  These physical-range counts are useful sanity checks but do not identify
the correct calibration; a sensor-specific Taupe/ARDP unit statement is still needed
before absolute residual weights are assigned.

Model entry should therefore be trend-oriented at first:

\begin{itemize}[leftmargin=1.5em]
  \item Compare simulated \(\Delta S_\ell\) or \(\Delta(nS_\ell)\) in EDZ bands with
  Taupe temporal trends, using \path{taupe_tdr_trend_operator.csv} as the current
  row-level operator artifact.
  \item Use closed-twin stability and open-twin seasonal amplitudes as qualitative or
  semi-quantitative validation.
  \item Do not force absolute Taupe values into the likelihood until the ARDP/Taupe
  workbook columns are mapped to a calibrated dielectric-to-water-content relation.
\end{itemize}

\subsection{Suction and relative humidity: pressure boundary and retention curve}

The suction/RH folder contains the open-twin sheets \texttt{OT\_RH5.xlsx} through
\texttt{OT\_RH8.xlsx}, a suction-location figure, the phase 30 report, TD slides, and
minutes.  The four RH files cover late 2021/early 2022 to 2025-09-04.  RH5 and RH6
stay near 93--96\%; RH7 and RH8 include low values near 13\%, which should be treated
as likely outliers or sensor/logger failures until confirmed.  The processed
\path{rh_open_twin_kelvin.csv} table contains 4,247 daily sensor rows, flags low-RH
outliers and values above the open-twin 95\% reliability caution threshold, and stores
the Kelvin-equation liquid-pressure conversion using \(T=298.15\,\mathrm{K}\) and
\(\rho_\ell=1095\,\mathrm{kg\,m^{-3}}\).

The current OGS model already uses a time-dependent pressure curve on the open niche.
The active include in \path{08_curves.xml} is
\path{08_08_open_niche_seasonal.xml}; the shifted
\path{open_niche_seasonal_curve_shifted.xml} is present but commented out in the
recovered projection setup.  The active curve contains 845 liquid-pressure values
ranging approximately from \(-5.23\) MPa to \(-53.24\) MPa.  Those values are
consistent in scale with a Kelvin-equation conversion from relative humidity to liquid
pressure/suction rather than with atmospheric pressure.  The Kelvin
relationship follows from vapour-pressure equilibrium over a curved liquid surface
\cite[pp.~448--452]{Thomson1871Kelvin}.  In Richards form, gas pressure is treated as
constant and \(\pcap=-p_l\); therefore an RH-derived suction can be applied as a liquid
pressure boundary \(p_l=p_g-p_c\), or as a negative gauge liquid pressure if \(p_g\) is
absorbed into the convention.  Ziefle et al. explicitly describe the CD-A model as
using Richards' equation, Van Genuchten retention, and Kelvin-equation boundary
conversion for desaturation around the niches \cite[pp.~4447--4448]{WaterContentEDZ2024}.

The working folder now audits this boundary condition explicitly in
\path{inversion_workflow/scripts/audit_rh_boundary_curve.py}.  Using the timestep
comment \(t_0=\) 2019-09-18, the script compares the active OGS curve with the
processed OT\_RH5--8 Kelvin pressures.  Of 4,247 RH rows, 2,280 lie inside the active
curve time range, 1,948 later rows are outside it, and 19 low-RH rows are excluded as
outliers.  Over the overlapping period the median absolute residual between
RH-derived gauge liquid pressure and the OGS curve is 13.00 MPa, with a mean absolute
residual of 16.32 MPa.  Therefore the current open-niche curve should not be claimed
as a verified reconstruction of the local RH workbooks; it is an existing model
boundary condition whose provenance or preprocessing needs confirmation from BGR/Gesa
before it is used as a fixed forcing in a calibrated inversion.

The follow-up semantics audit \path{build_rh_semantics_audit.py} writes
\path{rh_measurement_semantics_row_audit.csv},
\path{rh_measurement_semantics_sensor_summary.csv},
\path{rh_boundary_curve_semantics.csv},
\path{rh_measurement_semantics_summary.json}, and
\path{rh_measurement_semantics.md}.  It records 4,228 valid non-low-outlier RH rows,
19 low-RH outlier rows, and 2,492 open-twin rows above the 95\% thermo-hygrometer
caution threshold.  Inverting the active OGS curve through the same Kelvin coefficient
gives an implied RH range of 70.23--96.59\%; 772 of the 845 active curve rows are
drier than the clean RH5/RH6 workbook minimum.  This makes the provenance problem more
specific: the active curve may include older, filtered, or different boundary
humidity data, but it is not a direct reconstruction of the current OT\_RH5--8 sheets.

The request package \path{build_rh_boundary_provenance_request.py} turns that
mismatch into an external information request.  It writes
\path{rh_boundary_provenance_request.csv},
\path{rh_boundary_provenance_evidence.csv},
\path{rh_boundary_provenance_request_summary.json}, and
\path{rh_boundary_provenance_request.md}, with catalogue copies in the suction/RH
\path{derived_files} folder.  The package has six request rows and ten evidence rows.
It asks for the active curve source data and generation script, time-axis and
extension policy, RH sensor screening, Kelvin temperature/density constants,
open/closed curve mapping, and the retention-calibration release gate.  The generated
summary dates the active curve from 2019-09-18 to 2023-12-26 in model time, while the
copied RH workbooks continue to 2025-09-04.

The local workflow now also writes candidate RH-derived boundary curves with
\path{build_rh_boundary_candidate_curves.py}.  These are not accepted replacement
forcings; they are reproducible alternatives that expose what the currently copied
sensors would imply.  Six sensor-policy curves are generated, with OGS-style XML
snippets in \path{rh_boundary_candidate_curve_xml}.  The policy-preferred candidate is
\path{rh5_rh6_median}, because RH5/RH6 are the cleanest copied open-twin
thermo-hygrometer streams.  It contains 1,063 daily rows from 2021-12-16 to
2025-09-04, 576 overlap rows against the active curve, 487 rows after the active
curve ends, and an overlap MAE of 15.15 MPa.  This reinforces the current gate: a
locally reproducible RH-derived forcing can be prepared, but it should not replace
\path{08_08_open_niche_seasonal.xml} or enter a retention likelihood until the active
curve provenance, time axis, sensor screening, conversion constants, and extension
policy are confirmed.

The candidate policies are also compared as an uncertainty envelope by
\path{build_rh_boundary_uncertainty_audit.py}.  It writes
\path{rh_boundary_uncertainty_envelope.csv},
\path{rh_boundary_uncertainty_audit.csv},
\path{rh_boundary_uncertainty_summary.json}, and
\path{rh_boundary_uncertainty.md}, again with catalogue copies in the suction/RH
\path{derived_files} folder.  Across 1,064 daily envelope rows, 577 overlap the
active OGS curve and 487 extend beyond it.  In the overlap, the local candidate
pressure envelope has median width 2.10 MPa and 90th-percentile width 2.20 MPa, while
the active curve falls outside that local RH-derived envelope on 575 of 577 dates.
The active curve's mean absolute mismatch to the candidate-envelope median is
15.22 MPa.  This is not a reason to tune a residual directly against RH yet; it is a
quantitative expression of the provenance gate and supports requesting the original
curve-generation data/script before releasing boundary or retention parameters.

For inversion, RH/suction data enter in two ways:

\begin{itemize}[leftmargin=1.5em]
  \item \textbf{Boundary forcing:} reconstruct or audit the open-niche pressure curve
  from measured RH and temperature.  This is not a calibrated parameter unless the RH
  measurement uncertainty is explicitly included.
  \item \textbf{Retention validation/calibration:} compare measured suction/RH states
  with OGS \(S_\ell(p_c)\) from the Van Genuchten curve.  The current XML uses
  residual liquid saturation \(S_{\ell r}=0.1\), residual gas saturation \(0\),
  exponent \(m=0.45\), and \(p_b=10\) MPa, consistent with the van Genuchten form
  \cite[pp.~892--895]{VanGenuchten1980}.
\end{itemize}

The 2026 slides state that thermo-hygrometers are reliable in the open twin below
95\% RH, while psychrometers are preferred in the closed twin above 95\% RH.  That
reliability split should be encoded in the measurement error model
\cite[PDF p.~6]{TDMinutes2026Local}.

\subsection{Coordinates, bedding, and structural constraints}

The coordinate folder provides the essential measurement-to-model map.  The September
coordinate workbook has 109 entries: 19 NMR, 8 suction, 8 Taupe, and 72 permeability
entries.  It contains original/3D coordinates plus two 2D model coordinate blocks.
Email context says the \texttt{2D\_Model (x/y)} columns are the relevant block for the
current HERMES modelling.  The November TeamBeam coordinate workbook connects NMR,
Taupe, and characterization borehole labels to BGR model coordinates.  The processed
\path{measurement_coordinates_xy.csv} and \path{borehole_coordinates.csv} tables keep
these coordinate rows in flat form for mesh lookup.  The working folder now also
contains \path{inversion_workflow/scripts/build_mesh_observation_lookup.py}, which
maps those rows to the projected OGS bulk mesh.  It writes point lookups for 107
catalogue coordinate rows, endpoint lookups for 35 NMR/characterization/Taupe labels,
and 0.1 m line samples for 8 borehole or Taupe segments.  The projection mesh is a
local \([-5,5]\times[-5,5]\,\mathrm{m}\) domain; in the current lookup, 50 of the 107
catalogue coordinate rows fall outside that mesh bounding box.  These outside flags
are not errors in the catalogue, but they prove that only a subset of the collected
coordinate information can be compared directly with the current Niche-4 OGS domain.

The bedding/geology folder provides the structural prior.  Gesa's email gives an
approximate bedding angle of \(144^\circ\); the characterization papers describe the
upper sandy facies, fault zones, bedding-parallel structures, and the scale dependence
of geophysical and permeability measurements.  The CD-A characterization paper states
that the niches are oriented perpendicular to bedding strike and that the measurement
campaign includes ERT, MSM, NMR/CCM, Taupe, pulse tests, suction, and deformation
monitoring \cite[PDF pp.~2, 4--5]{Ziefle2024Characterization}.  For inversion this
means the anisotropy angle should be bedding-informed, while local anomalies near
faults should be either represented as heterogeneous permeability features or assigned
larger observation uncertainty.

\subsection{Projection workflow and mesh-based parameter fields}

The folder \path{model_projection_inputs/source_files} is not a measurement folder in
the physical sense, but it is the most concrete implementation bridge between the
measurement catalogue and the frozen OGS process.  The projection archive has also
been copied into the report working folder at
\path{GESA_model_original/projection_on_mesh_2025-09-05}.  Its \path{README.txt}
states that the project file was changed to call a projected mesh, that
\path{03_parameters_TRM.xml} contains the function/projection specification, and that
\path{bulk_all.vtu} domain identification must be re-done against the projected mesh
with \path{identifySubdomains}.

The script \path{generate_projections.py} reads \path{bulk.vtu} with \path{meshio} and
writes \path{bulk_w_projections.vtu} \cite[PDF p.~20]{CDAModellingSlides2025}.  It creates cell-data arrays called
\path{n_rd} and \path{k_i_rd} on the triangle elements, using the mesh
\path{MaterialIDs}.  The corresponding \path{03_parameters_TRM.xml} switches
\path{phi} and \path{k_i} from constants to OGS \path{MeshElement} parameters with
field names \path{n_rd} and \path{k_i_rd}; the project file then loads
\path{bulk_w_projections.vtu} as the bulk mesh.  This is consistent with the OGS
parameter mechanism: official OGS documentation states that \path{MeshElement}
parameters are defined by cell-label data arrays in a loaded mesh, selected through
\path{field_name} \cite[Sec.~``MeshElement and MeshNode'']{OpenGeoSysParametersDocs}.

As a reproducible replacement for the random scalar example, the working folder now
contains \path{inversion_workflow/scripts/generate_anisotropic_permeability_field.py}.
That script writes a four-component cell-data tensor
\[
\texttt{k\_i\_rd}=(k_{xx},\,k_{xy},\,k_{yx},\,k_{yy})
\]
using the same component order as the constant XML tensor
\(\texttt{1.00E-19 0.00 0.00 0.40E-19}\).  It parameterises the field by a
lognormal geometric-mean permeability, a global anisotropy ratio, and a bedding angle.
The example configuration uses \(\theta=144^\circ\), \(r=2.5\), and
\(k_{\rm ref}=6.32\times 10^{-20}\,\mathrm{m^2}\), matching the current XML reference
values and Gesa's bedding-angle note.

For the requested inversion this archive should be treated as a prototype, not as the
final parameter generator:

\begin{itemize}[leftmargin=1.5em]
  \item It proves that heterogeneous fields can be introduced without changing the TRM
  process equations: replace the constant XML parameter by a mesh property with the
  same physical meaning.
  \item The recovered random \path{k_i_rd} field is scalar, isotropic, uncorrelated,
  and sampled in linear space.  The new workflow script replaces it with a log-space,
  spatially correlated tensor field; an actual inversion driver should still replace
  the prior/proposal draw by sampled or optimized fields conditioned on the data.
  \item The script has no fixed random seed, so it is not reproducible as written.
  \item The porosity-field example is useful only if porosity is later released as an
  unknown.  For the first permeability-focused inversion it should remain fixed unless
  the data prove permeability alone is not identifiable.
  \item Before running the projected model, \path{bulk_all.vtu} and all boundary
  subdomains must be checked after projection; otherwise boundary conditions can be
  applied to stale or mismatched submeshes.
\end{itemize}

The workflow now also prepares an OGS-ready run from the direct permeability diagnostic
field in \path{inversion_workflow/runs/direct_fit_observation_run}.  This run copy
loads the fitted \path{bulk_w_projections.vtu} and activates output of
\path{pressure}, \path{saturation}, \path{temperature}, \path{displacement}, and
\path{porosity}, which are the minimum fields needed for later RH/suction, ERT/NMR,
Taupe, and HM observation operators.  The original GESA project files remain
unchanged; the output-variable edit is confined to the run directory and changes only
what OGS writes.  The run-input audit
\path{inversion_workflow/scripts/audit_ogs_run_inputs.py} confirms that the project
mesh files exist, boundary mesh references resolve, the required output variables are
requested, and \path{bulk_w_projections.vtu} contains both \path{k_i_rd} and
\path{n_rd}.  The enhanced audit also checks every VTU XML header for point and cell
counts plus declared \path{DataArray} entries, verifies that the copied
boundary/support submeshes expose \path{bulk_node_ids} and \path{bulk_element_ids},
and records that the Dockerized-Apptainer OGS run accepted these project inputs with
return code 0.  It still flags a reproducibility warning: \path{bulk_all.vtu} and the
copied boundary/support VTUs are present but fail local \path{meshio} decoding because
of appended base64 padding, so downstream Python tooling should not depend on
\path{meshio} reads of those submeshes unless they are regenerated with
\path{identifySubdomains}.  The local environment does not currently contain an \path{ogs}
host executable: \path{audit_ogs_environment.py} searched \path{PATH}, \path{/usr/bin},
\path{/usr/local/bin}, \path{/opt}, and \path{/home/ber0061} to depth 5.  The same
audit now records \path{ogs_container_found_runtime_available} because it finds the collected
container image
\path{cda_knowledge_base/file_transfers/collected/2025-04-04_2d_model_container/apptainer_ogs6.5.4.sif},
whose embedded label identifies \texttt{ogs/ogs@6.5.4}
(\texttt{cb5b3235101edecf3ba55e1039fe3c19bc13c636}) and
\path{/usr/local/ogs/build/bin} on the container path.  No Apptainer/Singularity or
\path{run-singularity} runtime is installed locally, but the preferred backend is now
the already used \path{docker_apptainer_sif} path: Docker runs the collected SIF
through \texttt{ghcr.io/apptainer/apptainer:latest}.  The script
\path{inversion_workflow/scripts/run_ogs_model.py} now records that
Dockerized-Apptainer command explicitly, including the workspace bind, numeric user
mapping, SIF path, and optional timeout.  The initial bounded 60-second startup run
with this backend parsed the recovered OGS 6.5.4 project, entered the coupled TRM
solve, and wrote seven partial \path{.vtu} files; it has since been superseded by
complete direct, regularized, and first adaptive candidate runs.  The post-run
sampler \path{inversion_workflow/scripts/sample_ogs_state_outputs.py} inventories the
full OGS outputs and samples 448 lookup rows per output time step: 107 catalogue
measurement points plus 341 borehole/Taupe line samples.  The sampled outputs contain
pressure, saturation, temperature, and displacement fields; porosity is still not
available as a direct OGS output field in these VTUs, so the sampler uses the fixed
support-mesh \path{n_rd} porosity fallback for the model-side proxy
\(\theta=\phi S_l\).  This activates 192 NMR rows in complete candidate runs, but
\(\theta\) must still be interpreted with the NMR bound/interlayer-water caveat
discussed above.

Operationally, the clean route is to generate one projected mesh per inversion sample
or per ensemble member, preserve the original May 2025 XML semantics, and change only
the parameter source from a constant to mesh-cell data.  Observation operators for ERT,
NMR, Taupe/TDR, permeability intervals, and RH/suction should then read simulation
outputs at the monthly timesteps documented in the projection archive rather than
trying to force all measurements onto the OGS mesh beforehand.

The working folder implements this route in
\path{inversion_workflow/scripts/prepare_ogs_run.py}.  The script creates a
self-contained run directory, copies the projection model, regenerates
\path{bulk_w_projections.vtu}, and writes \path{RUN_MANIFEST.json}.  In the current
smoke test the generated bulk mesh contains \(\texttt{k\_i\_rd}\) as a four-component
tensor field on all 10,239 \texttt{triangle6} cells and \(\texttt{n\_rd}\) as a fixed
porosity field, so the projection XML has both mesh fields it reads.

For the non-permeability data, \path{build_state_observation_targets.py} now writes
\path{state_observation_targets.csv} and \path{state_observation_samples.csv}.  These
tables do not force all measurements into the same residual form.  They record NMR as
point \(\theta=\phi S_l\) targets, RH as Kelvin-pressure boundary forcing, Taupe/TDR as
baseline-normalized EDZ-band trend targets with absolute calibration pending, and ERT
as external resistivity-field comparisons with calibration, spatial lookup, and a
direct-run diagnostic ready under explicit transform/support assumptions.

The next workflow layer is now active for NMR where the OGS run covers usable
observation times and mapped support.  The completed Dockerized-Apptainer runs
produce 83 output VTUs and 37,184 sampled state rows per executed candidate.  The OGS
VTUs contain pressure, saturation, temperature, and displacement but not an explicit
porosity field, so \path{sample_ogs_state_outputs.py} uses the fixed support-mesh
cell field \texttt{n\_rd} from \path{bulk_w_projections.vtu} as the porosity fallback.
The script \path{evaluate_state_observation_targets.py} writes 11,490 evaluation rows;
for the current best lower-support loop candidate, 192 NMR rows are finite objective residuals
against \(\theta=\phi S_l\), while 4,247 RH rows remain boundary forcing, 1,691 ERT
rows remain external diagnostic targets pending transform/support/uncertainty
confirmation, 5,265 Taupe/TDR or unsupported rows are not usable for the current state
fit, and 95 NMR rows are outside the output time tolerance.
\path{assemble_inversion_objective.py} then combines this state layer with the direct
permeability fit.  For the best executed lower-support loop candidate, the combined
objective has two active components: 75 direct pulse-test rows with effective
duplicate-aware weight 52 and objective value 269.84 for \(\sigma=0.5\) log10
permeability units, plus 192 sampled NMR state rows with objective value 2886.54.
The total active objective is 3156.37.

The model-entry coverage audit in
\path{inversion_workflow/measurement_operator_coverage.md} makes this status explicit
across all collected measurement streams.  It lists nine observation groups.  For the
current best executed lower-support loop candidate, two groups have active objective rows:
the direct permeability pulse tests and sampled NMR state targets.  ERT is inventoried
and matched to field files and has theta-to-resistivity calibration, a spatial lookup,
and direct-run plus 66-run cross-candidate log-resistivity diagnostics, but still
needs transform/support/uncertainty confirmation before assigning residual weights;
Taupe/TDR is mapped to EDZ-band line samples and has a
baseline-normalized trend operator with a run-local A3/A4 comparison diagnostic, but
remains outside the active objective until the absolute unit/calibration is confirmed;
RH has been converted to Kelvin pressure, audited as boundary forcing, and expanded
into local candidate boundary curves, but it remains outside the active objective
until the active-curve provenance and extension policy are confirmed.  Coordinates,
bedding/geology, and the model-projection inputs are support layers.  The remaining HM
monitoring streams now have a structured layout/qualitative layer and levelling summary
rows, but they remain outside the active objective until the Geoscope time-series
exports and laser-scan statistical files are located.

The generated model-entry matrix in
\path{inversion_workflow/measurement_model_entry_matrix.md} joins this coverage map
with the residual/likelihood semantics, stream activation gates, and final
include/exclude register.  It keeps the nine collected measurement classes on one
row each: two classes currently contribute 267 active objective rows, three classes
are diagnostic or boundary-audit only, three are support/prior layers, and other HM
monitoring is not ready for hard residual weights.  The matrix is therefore the
report-facing index for deciding whether a measurement is active, diagnostic,
support-only, boundary-only, or awaiting a collaborator/provenance decision.

The separate likelihood/activation audit in
\path{inversion_workflow/measurement_likelihood_model.md} records the stricter
statistical decisions that should not be inferred from data coverage alone.  It has
seven stream-level rows and two active streams: direct permeability pulse tests and
sampled NMR, for 267 active objective rows.  The active permeability residual is
\(\log_{10} k_{\rm pred}-\log_{10} k_{\rm obs}\) with the current broad
\(\sigma=0.5\) log10 scale and duplicate-aware weights.  The active NMR residual is
still caveated: the bound-water sensitivity audit rules out
a naive absolute \(\theta=nS_\ell\) residual, and the 66-run bias/anomaly audit shows
that a bias-safe NMR diagnostic changes the top candidate.  NMR therefore requires
either a trend/anomaly residual or explicit bound/interlayer-water bias or model-error term;
ERT and Taupe/TDR are kept as forward
observation operators until support/calibration choices are confirmed; RH remains a
boundary-condition provenance audit rather than a cell-wise residual.

A dedicated stream gate audit now makes these decisions machine-checkable in
\path{inversion_workflow/measurement_stream_activation_gate_audit.md}.  It records
eight stream/support layers, 20 gate checks, and seven required failed gates.  The
only streams with active objective rows are direct permeability and NMR, both with
tracked caveats.  ERT and Taupe/TDR are diagnostic-only, RH is boundary-audit-only,
and other HM monitoring is not ready for hard residuals because no Geoscope,
laser-scan, or full levelling numeric exports with epochs, units, support, and
uncertainty are present.  This audit is the current blocker list for promoting the
diagnostic streams into a weighted all-measurement objective.

The blocker list has been converted into a reproducible closure package in
\path{inversion_workflow/measurement_gate_closure_request.md}, with the machine table
\path{inversion_workflow/measurement_gate_closure_request.csv} and an email-ready draft
in \path{inversion_workflow/measurement_gate_closure_email_draft.md}.  It contains
13 closure items: seven external BGR/provider requests and six internal modelling
decisions.  The high-priority requests cover the ERT coordinate transform/support
mask and covariance model, Taupe workbook unit/calibration, RH active-curve
provenance and uncertainty policy, Geoscope/laser/levelling numeric exports and
metadata, and the final NMR bound/interlayer-water residual definition.  The
internal decision register also records the separate NMR not-promoted-default
policy for the current report state.
The medium-priority items record the permeability gas-pulse model-error policy, optional
historical permeability endpoint geometry, and Taupe grouped weighting/support
decisions.  Until these items are closed, the report should describe the current
fit as a direct-permeability plus conditional-NMR inversion with ERT, Taupe/TDR,
RH/suction, and other-HM streams used as diagnostics or provenance checks.

Before treating those rows as external-only, the same gate criteria were applied to
the local source catalogue and to raw \path{/home/ber0061/Downloads} files.  The
catalogue recovery audit in
\path{inversion_workflow/local_gate_recovery_audit.md} rescanned 2,540 source/index
documents and found zero possible gate-closing evidence rows.  The Downloads recovery
audit in \path{inversion_workflow/download_gate_recovery_audit.md} scanned 533
CD-A/HERMES-relevant filename, archive-member, Office-XML, and text documents and
also found zero possible gate-closing evidence rows.  Its companion inventory
\path{inversion_workflow/download_gate_recovery_inventory.csv} records 25
SHA1-verified duplicate/catalogue rows, including
\path{ERT_meas_Niche_open (1).zip}, \path{ERT_meas_Niche_open (2).zip}, and
\path{003_Nov_2025 (1).zip}, plus the extracted/run-output directory
\path{/home/ber0061/Downloads/CDA_N4_2D_250403}.  That directory contains run-output
VTUs and duplicated base model files rather than provider answers to the transform,
uncertainty, calibration, provenance, numeric-export, or endpoint-geometry gates.
Thus no currently known local download changes the activation status of the seven
external gates.

The seven external rows are now split into five recipient-specific drafts in
\path{inversion_workflow/external_gate_request_pack.md}, with suggested contact
routes resolved from the local mail index.  All five drafts have a suggested
\path{To} route through \path{Gesa.Ziefle@bgr.de}; the ERT draft additionally lists
\path{Markus.Furche@bgr.de} as suggested \path{Cc}, while the remaining provider
routes keep explicit forwarding caveats because no more specific sender was
resolved locally.  The response-intake tracker
in \path{inversion_workflow/external_gate_response_intake.md} creates stream-local
provider-response folders and records one response-note template, acceptance test,
and refresh-command list per external request.  The dispatch-readiness audit in
\path{inversion_workflow/external_gate_dispatch_audit.md} checks those drafts against
the intake tracker before sending: all seven requests are ready across the five
drafts, all seven rows carry suggested \path{To} routes, zero dispatch checks fail,
and the Gmail draft tracker
\path{inversion_workflow/external_gate_gmail_drafts.csv} records five Gmail drafts
covering all seven request rows.  No request is recorded as sent and all seven
responses remain missing.
A point-in-time Gmail live-state audit in
\path{inversion_workflow/gmail_gate_live_state_audit.md} then checked these same
threads plus the separate CTE confirmation draft.  At 2026-06-01 03:55:46 CEST, all six
expected Gmail drafts were still present as drafts; a sent-subject search found zero
sent copies; the recent named-provider search found zero replies from the Gesa,
Markus, or Tuanny routes; and the recent CD-A/HERMES/TeamBeam context search returned
only the generated drafts.  This confirms the current response-intake status but does
not replace sending or answering the requests.

The consolidated send/response state is now captured in
\path{inversion_workflow/external_blocker_dashboard.md}.  It joins the external
request, response-intake, dispatch, Gmail live-state, local-recovery, Downloads
recovery, and CTE confirmation trackers into eight open blocker rows: seven external
measurement gates and one CTE confirmation row.  All eight rows remain unsent and
missing-response blockers; the dashboard therefore provides the current worklist,
not closure evidence.

The current executed-candidate evidence is also joined in
\path{inversion_workflow/cross_stream_candidate_scorecard.md}.  The scorecard
compares 66 common run ids across the active objective, the NMR bias/anomaly
diagnostic, ERT log-resistivity MAE, and Taupe/TDR trend MAE without changing the
likelihood.  It finds no candidate in the top 10 of every stream.  The active
incumbent
\path{local_basis_sampler_002_basis_024_det_l_0p0075_s_1p000} ranks first in the
active objective but only 14th in the NMR bias/anomaly diagnostic, 8th in ERT, and
29th in Taupe/TDR.  The best mean-rank compromise is
\path{local_basis_sampler_003_basis_029_det_l_0p0100_s_1p000}, with mean rank 10.4
and worst rank 18.  This is the present evidence that the active-objective incumbent
should remain conditional on the unresolved diagnostic gates rather than being
reported as the final all-measurement permeability field.

The same pause point is now checked by
\path{inversion_workflow/runs/cross_stream_hybrid_field_plan/CROSS_STREAM_HYBRID_FIELD_PLAN.md}.
That screen builds 15 log-magnitude hybrids from five cross-stream winner runs and
the active incumbent.  It derives magnitude from the actual \path{k_i_rd} tensor,
not from the stale \path{k_mag_rd} metadata field, and then preserves the active
tensor orientation and anisotropy ratio by scalar multiplication.  The best hybrid,
\path{cross_hybrid_mean_rank_all_streams_plus2_a0p25}, has direct objective
269.8180571059851, i.e. zero improvement over the active direct objective at the
reported precision.  Thus the cross-stream hybrids are useful as deliberate
diagnostic probes, but they do not justify another OGS batch for the active
direct-permeability plus raw-NMR objective alone.

The parameter-release audit in
\path{inversion_workflow/inversion_parameter_release_plan.md} then connects those
measurement decisions back to the actual OGS XML.  It writes 14 rows from the base
\path{03_parameters_TRM.xml}, media files, projection \path{MeshElement} parameter
file, and selected regularized candidate run.  The result is deliberately restrictive:
the only stage-1 active unknown is the \path{k_i_rd} tensor-magnitude field; the
  \path{n_rd} porosity field is a fixed support field even though sampled NMR state
residuals are active; tensor shape, van Genuchten \(p_b\), and exponent stay gated
until the final NMR policy, support decisions, and RH boundary-curve provenance are
settled; elasticity, Biot, and swelling remain later mechanical validation
candidates; and the open-niche pressure curve plus the suspicious
\path{CTE} value are confirmation/provenance blockers, not ordinary calibration
parameters.  The same audit now records two active objective components with
objective 3156.353066948979, including 75 direct permeability rows and 192 sampled
NMR rows from 83 OGS output times, and it blocks further same-support active-objective
OGS batches until support, likelihood, bounds, or stream-gate evidence changes.
The dedicated \path{inversion_workflow/cte_confirmation_request.md} package now
turns the \path{CTE} blocker into a provider-facing question routed to Gesa with
Tuanny as the model-setup contact in copy.  It asks whether
\path{CTE=1254.74} is an intended thermal-expansivity value in \(1/K\), a copied
\path{c_p_s} heat-capacity value, inactive in the intended run, or another
unit/convention.  Until the response is recorded, the frozen XML value is kept
unchanged and uninterpreted.

The same staging decision is now enforced on prepared runs by
\path{inversion_workflow/scripts/audit_inversion_release_gates.py}.  The audit is
called inside the candidate driver and also writes a global summary in
\path{inversion_workflow/inversion_release_gate_audit.md}.  For the current three
regularized candidate directories, first three adaptive candidate directories, two
local-refinement candidate directories, three local-bracketing candidate directories,
three optimizer-proposed candidate directories, three continuous-proposed candidate
directories, three lower-support continuous candidate directories, six
lower-support loop candidate directories, six broad continuous candidate
directories, three local-basis candidate directories, and thirty production sampler
candidate directories it passes 1300 checks:
\path{k_i} is the active
\path{MeshElement} field \path{k_i_rd}, \path{phi} remains fixed support
\path{n_rd = 0.105}, media and retention XML files match the projection reference,
and no blocked or later-release parameter has been silently changed.
The consolidated guardrail audit in
\path{inversion_workflow/frozen_model_measurement_inclusion_audit.md} then links
that release control back to the measurement catalogue.  It checks 75
\path{RUN_MANIFEST.json} files, the 65 release-gated OGS candidate runs, the current
positive-definite permeability-field package, the 206-file measurement-info mirror,
1,880 indexed archive members, and 34 workbook-sheet summaries.  The audit passes,
but it also records that the current field remains an active-objective incumbent
rather than a final all-measurement inversion because the ERT, Taupe/TDR, RH,
other-HM, endpoint-geometry, and \path{CTE} gates remain open.

Finally, the full objective is tracked by
\path{inversion_workflow/objective_readiness_audit.md}, generated with
\path{build_objective_readiness_audit.py}.  This audit is intentionally stricter than
a progress note: it maps the original report/model/inversion request to 13 evidence
rows.  At the present state it reports \path{completion_state = not_complete}.  The
clean report build, source-model recovery, measurement inventory, processed
observation tables, source coverage, release plan, and release gate are in place.
The recovered OGS 6.5.4 SIF has now completed direct, regularized, first adaptive,
local-refinement, local-bracketing, and optimizer-proposed candidate runs through
continuous, lower-support continuous, repeatable lower-support loop, and broad
continuous candidate runs
through Dockerized Apptainer, and sampled NMR rows are active in the combined
objective.  The full inversion remains
incomplete because the current comparison is still proposal/sampler-batch selection
rather than a converged posterior or optimizer trace, ERT/Taupe/RH remain gated, and
the diagnostic streams have not been calibrated into a common likelihood.
The remaining report and modelling questions are also consolidated in
\path{inversion_workflow/open_question_resolution_matrix.md}, generated with
\path{build_open_question_resolution_matrix.py}.  That matrix links the original
open-question list, source/citation checks, report-comment audit, external gate
trackers, and final-promotion criteria.  It records 23 rows: 6 locally
resolved/current-scope rows, 9 external send/response rows, 6 internal
policy/final-decision rows, one tracked current-policy caveat, one deferred
time-dependency row, and an explicit no-new-evidence closeout decision row.  It
therefore turns the remaining questions into provider-response, modelling-policy,
or conservative-closeout gates rather than unresolved report comments.
The local-basis sampler is now included as executed combined evidence: its selected
three-candidate batch passed release gates and shifted the active best combined
objective to 3156.35307.
The current field is packaged separately in
\path{inversion_workflow/current_permeability_field/CURRENT_PERMEABILITY_FIELD.md},
with the copied \path{bulk_w_projections.vtu} mesh, summary JSON files,
cell-wise statistics for the actual \path{k_i_rd} tensor, detailed direct and sampled
NMR residual tables, and a run-local \path{ogs_run_inputs/} snapshot of the exact
project file, XML includes, meshes, and curves used for the executed incumbent.  The
accompanying \path{packaged_file_manifest.csv} records SHA256 checksums for the
packaged field, audits, residual tables, summaries, and OGS input snapshot.  The
generated \path{CURRENT_FIELD_REPRODUCIBILITY_AUDIT.md} verifies 19 package,
manifest, snapshot, objective, release-gate, execution, and field checks with zero
required failures.  The package records 10,239 positive-definite
\texttt{triangle6} tensor cells, fixed anisotropy ratio 2.5,
bedding angle \(144^\circ\), fixed support porosity \(n=0.105\), and the same active
objective value.  It is therefore the present deliverable candidate field for
inspection and reruns, but still not a final all-measurement inversion field.
The visual inspection artifact in
\path{inversion_workflow/current_permeability_field/visual_inspection/CURRENT_FIELD_VISUAL_INSPECTION.md}
renders six cell-wise PNG maps from that packaged VTU, including log10
geometric-mean permeability, minor and major eigen-permeability, local-basis
increment, nearest-anchor distance, and local-basis weight support.  These maps are
used as QA for the active-objective field and do not change the promotion status.
The promotion decision is made machine-checkable in
\path{inversion_workflow/current_field_selection_audit.md}: that audit accepts the
packaged field as the current active-objective incumbent, but refuses promotion to a
final all-measurement field because the provisional NMR trend/anomaly residual and
the cross-stream mean-rank scorecard prefer different runs, while ERT, Taupe/TDR,
RH, other-HM, and external-response criteria remain gated.
The companion scenario audit
\path{inversion_workflow/conditional_field_selection_scenarios.md} enumerates eight
gate/outcome scenarios; the packaged field wins only the current raw-NMR active
objective, while five different fields win across the scenarios.  This makes the
non-final status a consequence of the measured stream conflicts, not just a
wording caveat.
Those five scenario-winning meshes are copied with common tensor statistics in
\path{inversion_workflow/conditional_field_candidates/CONDITIONAL_FIELD_CANDIDATES.md}
so the current active-objective field, promoted-NMR field, ERT-screen field,
Taupe-screen field, and mixed diagnostic field alternatives can be inspected without
modifying the exchanged GESA model.  The same package now writes per-candidate
\path{diagnostic_metric_evidence.csv} files and a combined
\path{conditional_field_candidate_metric_evidence.csv}: 25 active-objective,
NMR-label-bias, NMR-trend/anomaly, ERT, and Taupe metric rows are extracted from
the global diagnostic audits, with no missing metric-evidence rows for the five
scenario winners.
The follow-up field-difference audit in
\path{inversion_workflow/conditional_field_candidates/CONDITIONAL_FIELD_DIFFERENCE_AUDIT.md}
compares the four non-current scenario winners against the current field cell by
cell using the geometric-mean \(\log_{10}\) permeability of \(\mathbf{k}_{i,\mathrm{rd}}\).
One alternative, \path{local_basis_sampler_003_basis_029_det_l_0p0100_s_1p000},
is essentially the same local-basis field at this scale: its mean absolute
difference is \(8.5\times10^{-6}\) log units and no cell differs by more than
0.05 log units.  The broader gate-dependent alternatives are local but material:
the largest mean absolute difference is 0.0377448 log units, with at most 241
of 10,239 cells exceeding 0.05 log units and 233 cells exceeding 0.10 log units.
This supports keeping the current field available for inspection while still
refusing to label it as an all-measurement inversion result.
The final-promotion gate is consolidated in
\path{inversion_workflow/final_inversion_promotion_checklist.md}.  That checklist
records 17 criteria: five pass, one passes with the tracked source-model caveat, and
two pass only for the active objective; six remain externally blocked, one remains an
internal NMR-residual policy decision, and two fail promotion because the scenario
winner is unstable and the current field-selection audit refuses final promotion.
Its current decision is therefore \path{do_not_promote_current_field}.  The checklist
is the practical close-out map: either the ERT, Taupe/TDR, RH, other-HM,
permeability-endpoint, CTE, and NMR-policy gates must be closed, or the modelling
team must explicitly record which gated streams are excluded from the final
objective.
The companion close-out playbook,
\path{inversion_workflow/final_inversion_closeout_playbook.md}, expands the nine
open promotion criteria into actions: six existing Gmail draft/response routes, one
internal NMR policy decision, and two scenario/final-field decision steps.  It lists
the response-note locations, acceptance evidence, and regeneration commands needed
after replies or explicit exclusions are recorded.  The regenerated playbook also
threads the direct-permeability support-conflict evidence through the permeability
endpoint, scenario-stability, and final-field approval rows, so resolving or
explicitly retaining the current support/likelihood policy is a named close-out
action before same-support OGS spending or final promotion can resume.
The additional decision register
\path{inversion_workflow/final_objective_decision_register.md} records the explicit
include/exclude layer for those same open rows.  It currently contains nine pending
or not-ready decisions: six external stream/provenance decisions, one internal NMR
policy decision, and two scenario/final-field decisions.  Each row states the
current allowed use of the evidence, the include path, the diagnostic-only or
exclusion path, the acceptance evidence, the report wording requirement, and the
refresh commands.  It now also carries the direct-permeability support-conflict
numbers and the still-unapproved support/likelihood policy into the endpoint,
scenario-stability, and final-field decision rows.  Consequently, a gated stream is
not silently waived: it must be
closed, kept diagnostic-only, explicitly excluded, or waived by a documented
modelling decision before the current permeability field can be promoted.
The companion final-objective scenario matrix
\path{inversion_workflow/final_objective_scenario_matrix.md} maps those
include/exclude choices to nine explicit selection options.  The current packaged
field wins only the narrow raw-NMR active-objective option and the direct-only
permeability tie case; every currently scored option that promotes NMR trend/anomaly
or accepts ERT or Taupe/TDR chooses a different run, while RH, other-HM, and
endpoint-missing historical permeability remain unscored future options until their
provider evidence is filed.  The matrix also records that the direct-only tie is
not a clean final objective: the current support map has repeated-row conflicts,
zero same-support reducible gap, and no approved direct-permeability policy change.
The recommendation packet
\path{inversion_workflow/final_objective_include_exclude_recommendations.md}
records the conservative no new evidence interpretation of the same rows.  It
recommends keeping ERT, Taupe/TDR, RH/suction, other-HM monitoring,
endpoint-missing historical permeability rows, and the suspicious CTE value out of
final hard likelihood weights unless provider evidence arrives or an explicit
modelling decision is recorded; it also keeps the current direct-permeability
support/likelihood policy as a signoff issue rather than a silent final assumption.
This packet is not itself a decision and does not
unblock promotion; it keeps the current field labelled as an active-objective
incumbent.
The further no-new-evidence closeout draft
\path{inversion_workflow/final_objective_no_new_evidence_closeout_draft.md}
translates that conservative position into exact decision text for review.  It
contains nine draft rows and a single-block acceptance text that would choose the
narrow \path{F01_current_raw_nmr_exclude_gated_streams} option only after
user/modelling-team approval and regenerated scenario, current-field, and promotion
audits.  The draft records no actual decisions, sends no mail, and does not promote
the current field.
The matching acceptance-record template
\path{inversion_workflow/final_objective_no_new_evidence_acceptance_record_template.md}
is the signoff guardrail for that path.  It has nine fillable rows, zero approvals
recorded, and a not-ready-to-apply status, so it prevents the draft text from being
mistaken for an accepted final objective.
The regenerated draft and acceptance template also carry the direct-permeability
support-conflict evidence into the final signoff layer itself: 75 active pulse-test
rows occupy 30 OGS support cells, 24 of those cells contain repeated rows, and 16
repeated cells span at least two \(\log_{10}\) permeability units.  The dominant
signoff example is still cell 4648 in BCD-A32 at 0.85--0.87 m, where three active
rows span 6.948847 \(\log_{10}\) units and include the configured-scalar envelope
conflict.  Because the permeability likelihood-policy acceptance template still has
0/1 primary approvals and the next field-fit gate marks same-support
active-objective batches as non-executable, even the conservative
no-new-evidence path cannot be read as approval to rerun the same support map or to
promote the active field.
For user review before any send action,
\path{inversion_workflow/gmail_draft_send_review_packet.md} consolidates those six
Gmail drafts into one table: five external measurement-gate drafts and one CTE
confirmation draft, eight covered request rows, all six marked ready for review, and
all six still recorded as unsent.
The production sampler/convergence audit now adds five executed production rounds
and recommends pausing the active direct-permeability/raw-NMR sampler.  The
cross-stream scorecard then shows why this pause is not a final inversion claim:
the active incumbent does not dominate the NMR bias/anomaly, ERT, and Taupe/TDR
diagnostics viewed together.
Two cheap follow-up field-family screens were therefore run before spending more
OGS time.  The cross-stream hybrid screen in
\path{inversion_workflow/runs/cross_stream_hybrid_field_plan/CROSS_STREAM_HYBRID_FIELD_PLAN.md}
blends the active field magnitude with diagnostic-winner magnitude patterns and
finds no direct-permeability improvement: the best hybrid ties the active direct
objective at 269.818057.  The structural/EDZ screen in
\path{inversion_workflow/runs/structural_edz_field_family_plan/STRUCTURAL_EDZ_FIELD_FAMILY_PLAN.md}
then tests 234 central EDZ caps, central shells, bedding-parallel bands at
\(144^\circ\), and broad BCD-A32/BCD-A33 open-twin corridors.  It catalogues 30
active permeability objective support cells and 1,928 ERT ready-support cells, but
again finds zero improving direct candidates; the best bedding-band row only ties
the active direct objective.  This makes the pause stronger: more OGS runs in the
current magnitude-field neighborhood are not justified by the active direct screen
alone, while ERT, Taupe/TDR, RH, and other-HM gates remain unresolved.
The residual conflict audit in
\path{inversion_workflow/permeability_residual_conflict_audit.md} explains why this
is not merely a sampler-convergence statement.  For the 75 active pulse-test rows it
finds weighted RMSE 1.610715 log units, 48 rows with absolute residual at least one
log unit, 21 rows at least two log units, and two high-permeability rows above the
configured scalar tensor range.  It also finds 24 support cells with repeated active
rows and 16 cells where repeated observations span at least one log unit.  Thus the
remaining direct mismatch is a measurement-support, gas-to-water interpretation,
uncertainty, or parameterization problem before it is another routine OGS batch.
The spatial support-conflict audit in
\path{inversion_workflow/permeability_support_conflict_spatial_audit.md} maps those
support cells back onto the current \path{triangle6} mesh.  It confirms that the 75
active rows occupy 30 support cells on the 10,239-cell current field, that 24 of the
support cells have repeated rows, and that 16 repeated cells span at least two
\(\log_{10}\) units.  The dominant mapped conflict is cell 4648 in BCD-A32 at
0.85--0.87 m, where three active rows span 6.948847 log units and include the only
configured-scalar envelope conflict.  This turns the stop rule into a spatially
inspectable statement: under the present one-value-per-support-cell map, the largest
direct residuals are co-located contradictory observations, not missing smooth-field
degrees of freedom.
The companion likelihood-policy audit in
\path{inversion_workflow/permeability_likelihood_policy_audit.md} makes this
distinction sharper without changing the active objective.  Under the recorded
duplicate-weighted Gaussian row policy the direct objective is 269.818057 and the
top ten rows contribute 0.494 of the row loss.  A support-cell weighted-mean
diagnostic, however, collapses the 75 active rows to 30 model-support groups and
has an objective-like value of \(1.77\times10^{-28}\), while a support-cell median
diagnostic remains 134.727790.  The current field is therefore matching the
duplicate-weighted means of conflicting rows that occupy the same support cells;
switching to robust tails, support-cell aggregation, or scalar-range outlier
handling must be an explicit likelihood decision rather than an implicit field-fit
success.
The support lower-bound audit in
\path{inversion_workflow/permeability_support_lower_bound_audit.md} turns this into a
field-search limit.  Holding the current observation-to-support-cell map fixed, the
best possible single log-permeability prediction for each support cell is the
duplicate-weighted observed mean.  The accepted field already attains that lower
bound: the current row-Gaussian direct objective and the single-support lower-bound
objective are both 269.818057, with zero same-support reducible gap.  All 30 support
groups are at their lower bound, and 16 groups still span at least two log10 units.
Thus another field that preserves the same support map cannot reduce the dominant
row-Gaussian pulse-test loss; reducing it requires changing support, gas-pulse
interpretation, likelihood semantics, scalar bounds, or tensor-shape freedom.
The decision request in
\path{inversion_workflow/permeability_likelihood_decision_request.md} keeps the
rowwise Gaussian policy as the current-report default, but records the required
alternatives before further active-objective OGS spending: robust row residuals,
support-cell aggregation, explicit scalar-range outlier handling, or a genuinely new
field-family gate with a materiality threshold.
The configured-scalar outlier disposition in
\path{inversion_workflow/permeability_configured_scalar_outlier_disposition.md}
then makes the scalar-bound issue local and explicit.  The two active rows outside
the configured scalar envelope are one duplicated BCD-A32 \(0.87\,\mathrm{m}\)
high-permeability value.  It is only \(0.107\) log10 units above the configured
upper envelope, whereas the same OGS support cell spans \(6.949\) log10 units
because a nearby active low-permeability row is mapped to the same cell.  The local
disposition is therefore not to widen eigenvalue bounds or release tensor shape
from these rows alone; they remain visible under the current rowwise Gaussian
default until the modelling team accepts that disposition or chooses a
robust/capped/support-cell/outlier policy explicitly.
The no-OGS existing-field rerank in
\path{inversion_workflow/permeability_likelihood_scenario_rerank.md} then scores
522 materialized permeability-field VTUs under those diagnostic policies.  It finds
40 active row-Gaussian best ties, with the current accepted field included in that
tie set.  Huber and Student-t row policies keep their representative winners inside
the tie set, whereas capped row loss, support-cell median aggregation, and
configured-scalar inside-only scoring select winners outside it.  Thus the rerank
is useful decision evidence, but not an approval to change the current-report
likelihood policy.
The winner cross-stream audit in
\path{inversion_workflow/permeability_likelihood_winner_cross_stream_audit.md}
then checks whether those policy winners have executed-run state and diagnostic
evidence.  Four of seven policy-winner rows have cross-stream scorecard evidence,
representing two unique fields.  All three outside-tie non-default winners are
direct-only fields without OGS/state/ERT/Taupe/NMR scorecard evidence.  The
row-Gaussian representative selected by candidate-id sorting has active-objective
rank 45 and mean all-stream rank 30.4, whereas the current accepted field has
active-objective rank 1 and mean all-stream rank 13.2.  A direct-only non-default
winner therefore cannot replace the current field without OGS execution and stream
diagnostics.
The next field-fit gate in
\path{inversion_workflow/permeability_next_field_fit_gate.md} makes this stop rule
explicit.  It records eight gates and recommends pausing same-support
active-objective OGS batches under the current rowwise Gaussian policy.  Further OGS
spending should be reopened only after a support mapping, likelihood, configured
bound/tensor-shape, or accepted measurement-stream objective changes; a non-default
direct-policy winner also needs OGS state and stream diagnostics before it can be
used for final promotion.
The consolidated packet
\path{inversion_workflow/permeability_likelihood_support_recommendations.md}
collects the same conclusion into eight direct-permeability decision-support rows.
It keeps \path{keep_rowwise_gaussian_default} as the current-report policy, records
zero same-support reducible gap, rejects immediate bound widening or tensor-shape
release, and states that the packet does not unblock final promotion.
The matching policy acceptance-record template
\path{inversion_workflow/permeability_likelihood_policy_acceptance_record_template.md}
is the signoff guardrail for any direct-permeability likelihood/support/outlier
policy change.  It contains the five exclusive policy options, currently records
zero of one required primary approvals, is not ready to apply, records no actual
decision, changes no active objective, and does not reopen same-support OGS
spending.

\subsection{Other HM monitoring}

The other HM monitoring folder contains minutes, TD modelling slides, levelling
slides, the HERMES input note, and the CD-A Tecplot layout.  These sources identify
extensometers, mini-piezometers, crackmeters, laser scans, levelling/miniprismas,
temperature/RH, and opening-time records.  The 2024 HERMES input note states that BGR
has long-term measurements of deformation, humidity, water content, pore water
pressures, crack width, pointwise piezometer/extensometer data, and two-dimensional
ERT/TDR/laser-scan data \cite[PDF p.~5]{InputHERMES2024Local}.

The working folder now makes this secondary HM layer explicit with
\path{build_other_hm_monitoring_inventory.py}.  The script streams
\path{VisualisationCDA.dat} and writes
\path{other_hm_visualisation_zones.csv}, \path{other_hm_visualisation_text_labels.csv},
\path{other_hm_levelling_displacements.csv},
\path{other_hm_qualitative_targets.csv}, and \path{other_hm_monitoring.md}.  The
Tecplot layout contains 84 zones and 11 legend labels.  The classified zones include
5 extensometer zones, 1 mini-piezometer zone, 22 fracture/crack-geometry zones, 12
miniprisma/geodetic zones, 2 laser-scan surface zones, 6 RH/suction-support zones, and
support zones for NMR, Taupe/TDR, permeability, convergence points, evapometer, and
niche geometry.

The 2026 TD minutes are important because they distinguish available streams from
usable residuals.  They state that a Geoscope update sent on 2026-04-20 covered RH,
temperature, door openings, extensometer, mini-piezometer, crackmeter, suction, and
laser scans with statistical interpretation \cite[PDF p.~3]{TDMinutes2026Local}.
Those raw Geoscope and laser-scan exports are not present in the current catalogue.
The request package
\path{inversion_workflow/processed_observations/other_hm_missing_numeric_request.md}
therefore lists the missing numeric exports explicitly: mini-piezometer time series
for BCD-A28 to BCD-A31, extensometer time series with BCD-A9/B10 failure metadata,
crackmeter time series, auxiliary Geoscope RH/temperature/door/suction logs if
distinct from the already catalogued RH files, the 2026-04-20 laser-scan statistical
interpretation, and a full precision-levelling survey table.  These rows remain at
zero active objective weight until the numeric files include units, support geometry,
reference conventions and quality/status flags.
The local absence claim is now checked by
\path{build_other_hm_numeric_source_audit.py}.  It scans the ten current other-HM
source files, the seven PDF members inside the 2026 presentation ZIP, extracted text,
the Tecplot layout support, and the levelling table.  All six missing-export request
classes have some local support geometry or extracted labels, but zero are
hard-residual-ready.  The ZIP contributes no spreadsheet/CSV-style numeric members;
the only numeric-extension source in this folder is \path{VisualisationCDA.dat},
which supplies support geometry rather than dated Geoscope time series.  The audit
therefore keeps mini-piezometer, extensometer, crackmeter, laser-scan, auxiliary
boundary-context, and full levelling residuals inactive until the requested exports
with units, timestamps or survey epochs, reference conventions, uncertainty, and
quality flags are delivered.
The same minutes report that mini-piezometers BCD-A28 to BCD-A31 were working well,
that BCD-A9/B10 closed-niche horizontal/vertical extensometers failed since September
2025, and that crackmeter data show one ongoing closed-niche trend plus open-niche
seasonality \cite[PDF p.~3]{TDMinutes2026Local}.  They also record the higher-level
trend split: the open twin shows ongoing trend in extensometer, suction and Taupe,
whereas the closed twin is mostly static in extensometer, pore-pressure and suction
measurements except for one crackmeter position with ongoing closure
\cite[PDF p.~7]{TDMinutes2026Local}.

The levelling slides provide the only pointwise numeric deformation summary currently
available in this secondary-HM bundle.  The extracted table has 12 height-change
values from the first 2022-08-29/30 measurement to the 2026-03 campaign.  The largest
reported changes are CDA-O1 \(-2.1\) mm settlement and CDA-C4 \(+1.3\) mm
heave/uplift; the reported detectable displacement is 0.1--0.2 mm at 95\% confidence,
depending on the point \cite[PDF pp.~5--6]{LevellingTD2026Local}.  These values can
become deformation-validation targets only after the survey reference frame and the
OGS displacement-output support are aligned.

The BGR 2026 modelling slides set the same priority for future HM observation
operators: compare with extensometer displacements/strains, mini-piezometer pore
pressures, and suction-derived pressure in the desaturated zone
\cite[PDF p.~2]{BGRModellingTD2026Local}.  They also warn that a simplified
perfect-ring 2D model already showed simulated strain discrepancies relative to the
extensometer measurements \cite[PDF p.~15]{BGRModellingTD2026Local}.  This is a useful
guard against accepting a hydraulic/permeability fit that is mechanically implausible.

For the current permeability-focused inversion these streams are secondary, but they
should be kept as validation gates:

\begin{itemize}[leftmargin=1.5em]
  \item Mini-piezometers validate pore-pressure trends after the Geoscope pressure
  series are located.
  \item Extensometers/laser scans validate displacement and strain fields, with the
  September 2025 closed-niche extensometer failure kept as a data-quality boundary.
  \item Crackmeters and convergence observations constrain whether simulated
  deformation trends are qualitatively plausible, especially the open/closed twin
  contrast.
  \item Levelling values can check the sign and millimetre-scale magnitude of vertical
  displacement after coordinate/reference-frame alignment.
  \item Temperature/RH/opening-time records explain boundary-condition disturbances.
\end{itemize}

\subsection{Recommended staged inversion workflow}

\begin{enumerate}[leftmargin=1.5em]
  \item Freeze the May 2025 OGS XML and April mesh/project structure as the base model.
  \item Validate the local data inventory with
  \path{inversion_workflow/scripts/validate_observation_manifest.py}; this guards
  against silently fitting against missing or structurally changed measurement files.
  \item Regenerate the normalized observation tables with
  \path{inversion_workflow/scripts/build_processed_observations.py} after any
  catalogue update, and use \path{processed_observations/README.md} as the table
  index.
  \item Build a geometry map from \path{coordinates_geometry_layout/source_files} or
  from the processed coordinate CSVs, then run
  \path{inversion_workflow/scripts/build_mesh_observation_lookup.py} to separate
  in-domain measurements from nearest-cell fallback or out-of-domain points.
  \item Build explicit non-permeability state targets with
  \path{inversion_workflow/scripts/build_state_observation_targets.py}.  The current
  layer contains 11,490 target rows: mapped NMR water-content targets, RH
  Kelvin-pressure boundary-forcing rows, Taupe/TDR baseline-normalized trend targets,
  and ERT external field-comparison rows.
  \item Quantify the NMR bound/interlayer-water caveat with
  \path{inversion_workflow/scripts/build_nmr_bound_water_sensitivity.py} before any
  absolute NMR water-content row is allowed into the numerical objective.
  \item Rebuild the likelihood/activation audit with
  \path{inversion_workflow/scripts/build_measurement_likelihood_model.py} so each
  measurement stream has an explicit residual transform, scale/weighting rule,
  bias/model-error caveat, and activation gate before it can enter the objective.
  \item Rebuild the XML-linked parameter-release plan with
  \path{inversion_workflow/scripts/build_inversion_parameter_release_plan.py}.  This
  keeps the inverse problem staged: active \path{k_i_rd} magnitude fields against the
  direct permeability and sampled NMR objective first, fixed \path{n_rd} support, no
  further same-support OGS spending without a gate change, and delayed release of
  tensor shape, porosity, retention, boundary, or mechanical parameters until the
  recorded gates are met.
  \item Audit prepared run directories with
  \path{inversion_workflow/scripts/audit_inversion_release_gates.py}, or use the same
  gate through \path{evaluate_inversion_candidate.py}.  A failed release gate should
  stop candidate ranking because the run no longer represents the documented
  permeability-only stage.
  \item Refresh the objective-level readiness audit with
  \path{inversion_workflow/scripts/build_objective_readiness_audit.py}.  This audit
  is the current completion check: it should remain \path{not_complete} until OGS
  outputs exist, state likelihoods are activated where justified, the five current
  external-gate Gmail drafts and the separate CTE Gmail draft are reviewed/sent and
  answered, and the combined inversion has been run.
  \item Re-run the local and Downloads gate-recovery audits with
  \path{inversion_workflow/scripts/build_local_gate_recovery_audit.py} and
  \path{inversion_workflow/scripts/build_download_gate_recovery_audit.py} whenever
  new TeamBeam, Gmail, Thunderbird, or raw download files are added.  A duplicate
  archive or OGS run-output folder is not a gate closure unless it contains the
  gate-specific transform, uncertainty, calibration, provenance, numeric-export, or
  endpoint-geometry evidence and is filed through the response-intake tracker.
  \item Rebuild or replace the projection workflow from
  \path{model_projection_inputs/source_files} so heterogeneous permeability fields are
  mesh-cell data with documented field names and reproducible generation.
  \item Prepare one OGS run directory per permeability-field sample with
  \path{inversion_workflow/scripts/prepare_ogs_run.py}; each run must carry a
  \path{RUN_MANIFEST.json} recording the generated or fitted field and the output
  variables required by observation operators.
  \item Execute the prepared run with \path{inversion_workflow/scripts/run_ogs_model.py}
  using either a host OGS executable, a host Apptainer/Singularity runtime, or the
  recorded Dockerized-Apptainer path for the collected OGS 6.5.4 SIF, and keep
  \path{OGS_EXECUTION_STATUS.json} with the run outputs.
  \item Sample OGS outputs with \path{sample_ogs_state_outputs.py}, evaluate the
  state target layer with \path{evaluate_state_observation_targets.py}, and assemble
  the direct permeability and state terms with \path{assemble_inversion_objective.py}.
  \item Use \path{evaluate_inversion_candidate.py} as the per-candidate harness.  The
  verified runs on the current direct-fit and smooth-search meshes prepare run
  copies, audit run-input mesh fields and submeshes, evaluate direct permeability,
  execute or record the OGS command, sample available state outputs, audit the RH
  boundary curve, evaluate state targets, and write the combined objective summary.
  \item Use \path{run_inversion_candidate_search.py} as the deterministic search loop
  around that harness.  The current dry search ranks four smooth candidates and keeps
  \path{length_0p025m} first with objective 269.97.  In dry mode this is still a
  direct-permeability-only ranking; in execute mode with the collected SIF through
  Dockerized Apptainer, the same loop ranks by the combined permeability plus
  state-observation objective and can be replaced by a Bayesian optimizer or ensemble
  sampler.
  \item Use \path{rank_regularized_permeability_candidates.py} to choose a small
  first OGS comparison set from the direct-only candidates.  At minimum this should
  include the direct-misfit winner and one damped regularized candidate, because the
  state observations are needed to decide whether the strongest permeability update
  is physically acceptable.
  \item Prepare that comparison set with
  \path{run_regularized_ogs_candidate_set.py}.  The current executed set contains the
  full-shift, 0.75-shift, and 0.5-shift 0.025 m fields, each with audited run inputs,
  83 OGS output VTUs, sampled NMR rows, and a combined objective summary.  The first
  adaptive broader-support batch has also been executed and is worse than the
  then-best 0.025 m field.  The first local-refinement batch improves the smooth-field best to
  \path{length_0p013m} with combined objective 3159.66.  The first local-bracketing
  batch around that field is worse than this best, with best combined objective
  3167.21.  The finite-candidate Bayesian proposal batch is also worse, with best
  combined objective 3172.37.  The first continuous proposal batch improves the
  active best to \path{length_0p006m_shift_0p992} with combined objective 3156.78.
  The lower-support continuous batch further improves the active best to
  \path{length_0p004m_shift_0p995} with combined objective 3156.38, and the first
  repeatable lower-support loop improves it again to
  \path{length_0p003m_shift_1p006} with combined objective 3156.37.  The second
  repeatable lower-support loop did not improve this incumbent, and the refreshed
  broad continuous evidence batches also did not improve it; the best broad-loop row
  is \path{length_0p010m} with combined objective 3157.70.  The next step is to
  treat \path{basis_024_det_l_0p0075_s_1p000} as the current incumbent evidence.
  Five production handoff rounds did not improve it; the local-anisotropy screen
  and cross-stream hybrid screen also found no direct-improving follow-up family.  Use
  \path{PRODUCTION_SAMPLER_DECISION.md} as the current decision to pause the active
  direct-permeability plus NMR production sampler and move effort to gated
  measurement streams or a new field family.
  \item Reproduce the current open-niche pressure curve from RH/T data, or document the
  residual mismatch if the source data do not fully reproduce it.
  \item Start with a heterogeneous anisotropic permeability field and fixed porosity,
  retention, elasticity, swelling, and thermal properties.
  \item Use pulse-test data as direct log-permeability constraints with large
  interval-scale errors and gas/slip caveats.
  \item Use ERT through the documented theta-to-resistivity operator after spatial
  projection to the ERT/common mesh exists; use NMR as a sparse water-content anchor
  with bound-water bias/error.
  \item Use Taupe/TDR through the documented trend operator and suction/RH as
  trend/consistency constraints until their calibration to model saturation is
  documented.
  \item Only after permeability is identifiable should porosity, \(m\), \(p_b\), or
  other fields be released as additional unknowns.
\end{enumerate}
